% =====================================================================
% Second adversarial-wave refinement: corrections and extensions to the
% post-adversarial synthesis spine.
% Author: Raeez Lorgat.  Chriss--Ginzburg register; subscripted
% $\kappa_\bullet$; no AI attribution; no bookkeeping vocabulary.
%
% This section supersedes specific passages of the first-wave spine
% where a second adversarial pass exposed scope-slips, category errors,
% attribution mistakes, or arithmetic overstatements.  The Wave-1
% ghost-theorems survive; a second layer of ghost-theorems is extracted
% here from the Wave-1 healed statements themselves.
% =====================================================================

\section{Second-pass refinement of the platonic synthesis}
\label{wn:sec:spine-second-pass}

\emph{The first-wave synthesis survived a second adversarial pass with
thirty-three corrections, eight promotions, and three new frontier
openings. The corrections fall into four categories: category slips
(where an invariant was mis-attached to the wrong object), mechanism
errors (where a true conclusion had a false proof), attribution
mistakes (where a theorem was cited to the wrong primary source or to a
non-existent paper), and arithmetic overstatements (where a bound was
loose or a sequence was mis-listed). Each correction is stated as a
retraction with its ghost-theorem extracted.}

\subsection{The K3 $\times$ E crystal as a six-invariant four-functor
square}
\label{wn:subsec:second-pass-six-invariants}

\begin{theorem}[Four-functor square for the $K3 \times E$ modular
characteristics]
\label{wn:thm:second-pass-four-functor-square}\ClaimStatusTheorem

The first-wave "four-value crystallisation $\{2, 3, 5, 24\}$" is in
fact an invariant of six numbers $\{0, 2, 3, 4, 5, 24\}$ distributed
across \emph{two Calabi--Yau categories} via a four-functor square:
\[
\begin{array}{c@{\qquad}c}
d = 2 & d = 3 \\
\hline
D^b\,\mathrm{Coh}(K3) & D^b\,\mathrm{Coh}(K3\times E) \\
\kcat^{\mathrm{K3\text{-}fibre}} = \chi(\cO_{K3}) = 2
 & \kcat^{\mathrm{total}} = \chi(\cO_{K3\times E}) = 0\\[1.5ex]
\Phi^{\mathrm{FA}}_2 \downarrow & \Phi^{\mathrm{FA}}_3 \downarrow \\[1.5ex]
\cH_{\Mukr}(K3) & \fg_{\Delta_5}\\[1.5ex]
\kch^{\mathsf{B}} = c_+(\Mukr(K3)) = 4
 & \kch^{\mathrm{hyp}} = \mathrm{rk}_{\mathrm{hyp}}(\fg_{\Delta_5}) = 3
\end{array}
\]
The values $\{5, 24\}$ are shared invariants across both categories:
$\kBKM(\Phi_1) = 5$ is the Borcherds weight of $\Delta_5$ attached to
the $d = 3$ Stage-$2$ output, and $\kfib(K3) = \mathrm{rk}(\Wttl(K3)) =
24$ is the Mukai-lattice fibre rank of the $K3$ factor used as a
$(\Sigma_2, C)$ datum in the $d = 3$ specialisation.

\emph{The two additional invariants that emerge under audit.} The
Calabi--Yau-$3$ total-space $\kcat^{\mathrm{total}}(K3 \times E) = 0$
is forced by K\"unneth-multiplicativity of $\chi(\cO_X)$ on products
with trivial K3 factor; the $\mathsf{B}$-row $\kch^{\mathsf{B}} = 4$ at
$d = 2$ is the five-archetype landmark witness
$\cH_{\Mukr}(K3) = \Phi_2(D^b\,\mathrm{Coh}(K3))^{\mathrm{Heis}}$ on
the Mukai lattice $\Wttl(K3) \simeq E_8(-1)^{\oplus 2} \oplus
U^{\oplus 3}$ of signature $(4, 20)$, producing the Beilinson--Drinfeld
Koszul-conductor value $K^{\kch}_{\mathsf{B}} = \varrho K = 8$.

\emph{Stage-$2$ unit shift.} The numerical identity
$\mathrm{rk}_{\mathrm{hyp}}(\fg_{\Delta_5}) = c_+(\Mukr(K3)) - 1 = 3$
is a conjectural \emph{unit shift} produced by the Stage-$2$
specialisation $\mathrm{Sp}^{\mathrm{ch}}_{K3, E}$ from the $d = 2$
Mukai-enhanced Heisenberg on $K3$ to the $d = 3$ BKM on
$K3 \times E$. The shift is explained, at the conjectural level, by a
transverse-cycle-class reduction of the positive-definite rank-$4$
Mukai witness to the rank-$3$ hyperbolic real-root subalgebra of
$\fg_{\Delta_5}$ under the integration-over-$\Sigma_2$ operation.
\end{theorem}

\begin{corollary}[Every K3 $\times$ E reading is a declared scope]
\label{wn:cor:second-pass-scope-discipline}
When the monograph displays the four-value crystal $\{2, 3, 5, 24\}$
(manifesto convention), the reader should understand:
$2 = \kcat$ at K3-fibre scope;
$3 = \kch(\fg_{\Delta_5}) = \mathrm{rk}_{\mathrm{hyp}}$ at the
$d = 3$ BKM-Cartan scope;
$5 = \kBKM(\Delta_5) = c_1(0)/2$ at the automorphic-weight scope;
$24 = \kfib(K3)$ at the Mukai-lattice-rank scope.
Two further invariants live at complementary scopes:
$\kcat^{\mathrm{total}}(K3\times E) = 0$ and
$\kch^{\mathsf{B}}(\cH_{\Mukr}(K3)) = 4$. Each of the six values is an
invariant of a distinct mathematical object; the monograph flags the
scope at every use.
\end{corollary}

\subsection{The one Borcherds ladder, three compatible lifts}
\label{wn:subsec:second-pass-one-ladder}

\begin{theorem}[Single CHL Borcherds-weight ladder, three mutually
compatible primary lifts]
\label{wn:thm:second-pass-single-ladder}\ClaimStatusTheorem
On the CHL slice $N \in \{1, 2, 3, 4, 6\}$, the modular
characteristic $\kBKM(\Phi_N) = c_N(0)/2$ takes values
\[
(c_N(0))_{N=1,2,3,4,6} \;=\; (10, 4, 2, 2, 2)
\qquad\Longrightarrow\qquad
\kBKM(\Phi_N) \;=\; (5, 2, 1, 1, 1)
\]
on the singly-twined Jacobi-input convention
$\phi^{(g_N)}_{0, 1} = \tfrac{1}{2}\, Z^{(g_N)}_{K3}$
(Eichler--Zagier $1985$; Cheng--Harrison--Paquette $2014$ Table~$4$).
This single ladder is produced by three mutually-compatible primary
lifts, all landing on the same paramodular Borcherds product
$\Phi_N = \Delta^{(N)}$:
\begin{itemize}
\item Borcherds multiplicative singular-theta lift
(Borcherds $1995$ Thm.~$13.3$), weight input $1 - b^+/2 = -3/2$ on the
vector-valued Weil representation attached to $\Lambda^{3,2}$;
\item Gritsenko additive paramodular lift
(Gritsenko $1999$ Thm.~$1.1$, using the \emph{index-$2$} Jacobi input
$\phi^{(g_N)}_{k(N), 2}$ with integer weight $k(N)$);
\item Gritsenko--Nikulin direct CHL-paramodular construction
(Gritsenko--Nikulin $1998$ Thm.~$2.1$).
\end{itemize}

\emph{Retraction of the earlier two-ladder framing.} The first-wave
spine stated a second ladder $(5, 4, 3, 2, 1)$ via "Gritsenko
additive-lift of weight-$k(N)$ index-$1$ Jacobi forms". This is
incorrect at three points: (i) the cuspidal Jacobi space
$J^{\mathrm{cusp}}_{0, 1} = \{0\}$ is empty (Eichler--Zagier $1985$
Thm.~$3.5$), so no weight-$0$ index-$1$ input exists; (ii) the
Gritsenko additive lift preserves weight rather than shifting it, so
no weight-raising ladder arises from index-$1$ inputs; (iii)
$\Delta_5 = \mathrm{Grit}(\phi_{5, 2})$ uses index-$2$ input
(Gritsenko $1999$ Thm.~$1.1$), not index-$1$ (Thm.~$1.2$). The correct
hidden structure is the three-lift compatibility stated above; the
two-ladder framing was an artefact of conflating the refined
Gritsenko--Nikulin Thm.~$2.1$ values with the Gritsenko Thm.~$1.1$
values.
\end{theorem}

\begin{theorem}[Boundary-extension $N \in \{5, 7, 8\}$, corrected
direction]
\label{wn:thm:second-pass-boundary-extension}\ClaimStatusTheorem

The non-CHL rows at $N \in \{5, 7, 8\}$ carry half-integer weights and
live on the metaplectic cover $\Mpr_4(\mathbb{A})$, classified by the
metaplectic cocycle $v^{\theta}_N \in H^1(\Gamma^{(2)}_N, \{\pm 1\})$.
The factorisation is:
\[
S_{k + 1/2}(\Gamma_0(4N)) \xrightarrow{\;\text{Niwa }1974\;}
S_{2k-1}(\Gamma_0(2N)) \xrightarrow{\;\text{additive theta}\;}
M_{k + 1/2}(\Gamma^{(2)}_N, v^{\theta}_N),
\]
with Niwa $1974$ providing the half-integer-to-integer Shimura
correspondence direction, and the additive theta lift using the Howe
pair $(\Mpr_4, \mathrm{O}(1, 1))$ inside $\Spr_8$ (\emph{not}
$(\Mpr_2, \mathrm{O}(3, 2))$ as the first-wave spine stated).

\emph{Why Niwa is the correct primary.} Niwa's theorem integrates a
cusp form $g$ of weight $k + 1/2$ on $\Gamma_0(4N)$ against the
Shimura theta kernel $\theta_{\mathrm{Shim}}(z, \tau)$ and produces
a cusp form on $\Gamma_0(2N)$ of weight $2k - 1$ --- this is the
\emph{direction} the metaplectic paramodular lift requires, starting
from a half-integer input and recovering an integer-weight elliptic
form for the subsequent additive theta lift. Shimura's original 1973
theorem goes in the opposite direction (integer to half-integer), so
quoting Shimura here would invert the arrow; Niwa established
precisely the needed half-to-full direction by explicit seesaw pair
calculation. The Howe pair $(\Mpr_4, \mathrm{O}(1, 1))$ inside
$\Spr_8$ is forced: $\Mpr_4$ is the double cover of $\mathrm{Sp}_4$
acting on the genus-$2$ Siegel upper half-space, and the complementary
partner in the seesaw must have signature matching the square-root
of $\det$ multiplier, which is $\mathrm{O}(1, 1)$ --- not
$\mathrm{O}(3, 2)$, whose Howe partner is $\Mpr_2$ and which lives
in a different seesaw tower.

\emph{Retraction of earlier mechanism.} The first-wave spine
factorised through the Weil correspondence
$\theta_{2, 3}\colon \Mpr_2 \to \mathrm{O}(3, 2)$ with Shimura $1973$
at input and Waldspurger $1980$ at output. Both endpoints were in the
wrong direction: Shimura $1973$ lifts $S_{k+1/2} \to S_{2k}$ (not
$S_1 \to S_{1/2}$), and Waldspurger $1980$ concerns
$(\Mpr_2, \mathrm{PGL}_2)$-Howe duality (not
$(\Mpr_2, \mathrm{O}(3, 2))$). The correct lift at half-integer
weight is Niwa $1974$ at the input side and
Gritsenko--Nikulin $2002$ / Gritsenko--Cl\'ery $2018$ additive-theta
construction at the output side.

The trifurcation $F^{\mathrm{total}} = F^{\mathrm{SK}} \sqcup
F^{\Mpr} \sqcup F_{\mathrm{exc}}$ across
$\{N \in \{1,2,3,4,6\}\} \sqcup \{N \in \{5,7,8\}\} \sqcup
\{\text{exceptional Monster}\}$ survives; only the internal mechanism
of the $F^{\Mpr}$ arm is rectified.
\end{theorem}

\subsection{The $\Phi_{12}$/$\Phi_{10}$/$\Delta_5$ relation: Wang--Williams
pullback rigidity}
\label{wn:subsec:second-pass-wang-williams}

\begin{theorem}[Wang--Williams pullback rigidity for $\Phi_{12}$]
\label{wn:thm:second-pass-wang-williams}\ClaimStatusTheorem

The Borcherds $\Phi_{12}$ is the \emph{unique} holomorphic singular-weight
Borcherds product on $\mathrm{II}_{26, 2}$ (Wang--Williams $2023$
Thm.~$3.1$). \emph{Why uniqueness holds.} Singular weight on a
signature-$(b^+, 2)$ lattice means Borcherds-weight $= b^+/2 - 1$,
the smallest weight at which the Borcherds lift can be holomorphic;
at this threshold, the Fourier coefficients of the input vector-valued
modular form are rigidly constrained (all but the unique lattice-vector
class of negative norm $2$ must vanish), and holomorphicity at the
$2U$-cusp pins the remaining coefficient to an explicit normalisation.
For $\mathrm{II}_{26, 2}$ these constraints have exactly one solution.
Every holomorphic Borcherds product of singular weight on
a lattice $2U \oplus L$ that is non-vanishing at the $2U$-cusp is a
pullback of $\Phi_{12}$, with $L$ a finite-index sublattice of the
Leech lattice $\LLeech$ (Wang--Williams $2023$ Thm.~$3.5$). \emph{The
pullback mechanism.} A primitive lattice embedding $2U \oplus L
\hookrightarrow \mathrm{II}_{26, 2}$ sits over a finite-index
inclusion $L \hookrightarrow \LLeech$; the Borcherds lift commutes
with primitive pullback under $2U \oplus L$-restriction; non-vanishing
at the $2U$-cusp forces the pulled-back form to inherit the
singular-weight rigidity of $\Phi_{12}$.

\emph{Relation to $\Phi_{10}$ and $\Delta_5$.} The Igusa $\Phi_{10}$
weight $10$ on $\mathrm{II}_{3, 2}$ and the paramodular $\Delta_5$
weight $5$ on the paramodular cover satisfy
$\Phi_{10} = \Delta_5^2$ (Gritsenko--Nikulin $1998$). The interaction
with $\Phi_{12}$ is a wall-crossing identity on a tubular
neighbourhood of the Humbert divisor $H_1 \subset \cA_2$:
\[
\Phi_{12}(Z) \,/\, \Phi_{10}(Z)\bigr|_{H_1\text{-tubular}}
\;=\; \Psi_{\mathrm{wt}\,2}(Z),
\]
where $\Psi_{\mathrm{wt}\,2}$ is a weight-$2$ Borcherds lift
parametrising the wall-crossing between the $\mathrm{II}_{26, 2}$ and
$\mathrm{II}_{3, 2}$ Grassmannian components.

\emph{Retraction of earlier hyperbolic-restriction framing.} The
first-wave spine stated "$\Phi_{12}$ on $\mathrm{II}_{26,2}$ restricts
along $\mathrm{II}_{2,2} \hookrightarrow \mathrm{II}_{25,1}$ to
$\Phi_{10} = \Delta_5^2$" (Borcherds $1998$ \S$14$ K\"unneth). Four
structural errors: (i) signature $(2, 2)$ cannot embed in $(25, 1)$;
(ii) Borcherds $1998$ \S$14$ has no K\"unneth formula (it is the
Shimura--Doi--Naganuma--Maass--Gritsenko correspondence section);
(iii) pullback of a Borcherds product preserves weight, so
$\Phi_{12}$ (weight $12$) cannot pullback to $\Phi_{10}$ (weight $10$);
(iv) $\mathrm{II}_{25, 1}$ is the Fake Monster root lattice whereas
$\mathrm{II}_{26, 2}$ is the Grassmannian-domain lattice — different
structural r\^oles, related by $\mathrm{II}_{26, 2} = \mathrm{II}_{25,
1} \oplus \mathrm{II}_{1, 1}$. The ghost-theorem is Wang--Williams
pullback rigidity above; the true arithmetic linking $\Phi_{12}$ with
$\Phi_{10} = \Delta_5^2$ is the Humbert-divisor wall-crossing identity,
not a K\"unneth restriction.
\end{theorem}

\subsection{Fulton--MacPherson: strict operad composition, not BM
convergence}
\label{wn:subsec:second-pass-FM-reason}

\begin{theorem}[Why Fulton--MacPherson is the correct operadic base]
\label{wn:thm:second-pass-FM-reason}\ClaimStatusTheorem

The $E_3^{\mathrm{hol}}$-operadic composition on $\mathbb{C}^3$
requires the Axelrod--Singer / Fulton--MacPherson compactification
$\FMcpt(n)$ for reasons of \emph{smooth-manifold-with-corners
compatibility with operadic composition}, not for Bochner--Martinelli
convergence. The structural trichotomy is:
\begin{itemize}
\item On the open $\mathrm{Conf}_n(\mathbb{C}^3)$ with
$P_{\mathrm{BM}}$, the Feynman amplitudes converge as numbers via
$U(3)$-equivariant angular average at the codim-$6$ 2-point diagonal
(Axelrod--Singer $1994$ Lem.~$5.7$), but the codim-$12$ 3-point
diagonal contributes a genuine power divergence
$d\rho/\rho^{4}$ that angular averaging does not resolve.
\item On the Costello--Gwilliam heat-kernel regularisation at scale
$L > 0$, amplitudes are homotopy-coherent as a prefactorisation algebra
in $\mathrm{Ch}_{\mathrm{Dolb}}$ with $L \to 0$ recovery.
\item On $\FMcpt(n)$, the manifold-with-corners blowup structure
admits a strict symmetric monoidal operadic composition
$\gamma\colon \FMcpt(p) \times \prod_i \FMcpt(q_i) \to \FMcpt(q_1 +
\cdots + q_p)$ as a smooth map of corners, underwriting strict
chain-level $E_3^{\mathrm{hol}}$-operadic action.
\end{itemize}
All three routes agree at Dolbeault cohomology
(Gwilliam--Williams $2021$ Thm.~$2.5.5$); the difference is at the
chain level between $L_\infty$-coherences.

\emph{Retraction of earlier mechanism.} The first-wave spine
attributed the necessity of $\FMcpt$ to "conditional-vs-absolute
convergence" of $P_{\mathrm{BM}}$. The power-counting is: at the
codim-$6$ diagonal, $P_{\mathrm{BM}} \sim \|z - w\|^{-5}$ (not
$\|z - w\|^{-6}$), volume $\rho^{5}\, d\rho$, product $d\rho$,
marginally log-divergent --- saved by $U(3)$-angular average. At the
codim-$12$ diagonal, three propagators $\|z - w\|^{-5}$ against volume
$\rho^{11} d\rho$ give $d\rho/\rho^{4}$, genuinely power-divergent.
$\FMcpt$ is required at the three-point configuration, not by the
propagator singularity alone.
\end{theorem}

\subsection{The Bardeen--Zumino category error}
\label{wn:subsec:second-pass-BZ}

\begin{theorem}[The two one-loop representatives are genuinely
different invariants]
\label{wn:thm:second-pass-BZ-category}\ClaimStatusCorrected

The quadratic-Casimir wave-function renormalisation
$Z^{(1)}_{\cA} = 1 - \hbar C_2(\fg)(4\pi)^{-3}\log(L/\varepsilon)$
and the cubic-Casimir BV anomaly
$\kanom(X, \fg) = \hbar A(\fg)\chi_{\mathrm{top}}(X)(2(4\pi)^3)^{-1}
\|\Omega_X\|^2$
(with $A(\fg) = d^{abc}d_{abc}/\dim\fg$) are invariants of
\emph{different cohomological objects}:
the wave-function counterterm is a ghost-number-$0$ class in
$H^0_{\mathrm{loc}}(\cE_{\hCS})$ absorbed by field redefinition;
the BV anomaly is a ghost-number-$+1$ class in
$H^1_{\mathrm{loc}}(\cE_{\hCS}[-1])$ obstructing QME solvability.
No local cochain can conjugate between them: they live in different
graded sectors of the BV complex and represent different physical
phenomena (UV running of the normalisation versus obstruction to
gauge-invariant regularisation).

\emph{The true content of the classical Bardeen--Zumino polynomial}
(Bardeen $1969$; Zumino $1983$ \emph{Nucl.~Phys.~B} $223$) is to relate
the two \emph{representatives} --- consistent and covariant --- of the
same cubic-Casimir triangle anomaly, both of which live in the
ghost-number-$+1$ sector:
\[
Q_{\mathrm{BRST}}(\mathrm{BZ}^{\mathrm{hol}}) \;=\;
\kanom^{\mathrm{cons}} - \kanom^{\mathrm{cov}},
\qquad
\kanom^{\mathrm{cov}}, \kanom^{\mathrm{cons}} \in
H^1_{\mathrm{loc}}(\cE_{\hCS}[-1]).
\]
This intra-cubic bridge is a genuine theorem (Costello $2013$; see
\texttt{chapters/theory/quantum\_chiral\_algebras.tex} Remark
\texttt{rem:plat-Z-vs-anomaly}).

\emph{Retraction of earlier cross-sector bridge.} The first-wave spine
asserted
$Q_{\mathrm{BRST}}(\mathrm{BZ}^{\mathrm{hol}}) =
\kanom^{\mathrm{cov}} - \kanom^{\mathrm{cons}}$
bridging the quadratic-Casimir and cubic-Casimir sectors. This is a
category error: the two sides have different ghost numbers and
different Feynman-graph topologies (two-leg bubble versus three-leg
wheel), and the BRST differential cannot conjugate between them. The
ghost theorem is the intra-cubic consistent-vs-covariant bridge stated
above.
\end{theorem}

\subsection{Formality on $\mathbb{C}^3$ and on $K3 \times E$: correct
sources, correct receptacle}
\label{wn:subsec:second-pass-formality}

\begin{theorem}[Flat-space formality and the $K3 \times E$ Yukawa]
\label{wn:thm:second-pass-formality}\ClaimStatusTheorem

\emph{Flat $\mathbb{C}^3$.} The minimal $L_\infty$-model of
$\cE_{\hCS}(\mathbb{C}^3) = \Omega^{0,\bullet}_c(\mathbb{C}^3,
\fg)[1]$ has $\ell_n^{\min} = 0$ for all $n \geq 3$ by a
Dolbeault-degree count on the Hodge-homotopy-regularised sector
(Kontsevich $1997$ q-alg/9709040 \S$6.4.1$ lifted to Dolbeault by
Costello--Gwilliam $2017$ Vol.~II Ch.~$7$). No canonical Hodge
homotopy exists on the full Fr\'echet space
$\Omega^{0,\bullet}(\mathbb{C}^3, \fg)$; the construction is performed
on the compactly-supported sector via BV-regularised heat-kernel
homotopy $h_{\varepsilon} = \bar\partial^{\ast}
\int_0^{\varepsilon} e^{-s\Box_{\bar\partial}}\, ds$
(Costello--Li $2016$ \S$3$), with $\varepsilon \to \infty$ limit on
polynomial harmonic representatives giving the Bochner--Martinelli
kernel.

\emph{Compact $K3 \times E$.} The cubic $L_\infty$-bracket receptacle
on the product decomposes via K\"unneth as
\[
H^2(K3, \Lambda^2 T_{K3}) \otimes H^1(E, T_E)
\;\simeq\; H^2(K3, \cO_{K3}) \otimes H^1(E, \cO_E)
\;\simeq\; \CC \otimes \CC
\;\simeq\; \CC,
\]
using $\Lambda^2 T_{K3} \simeq \cO_{K3}$ (K3 is CY$_2$) and
$T_E \simeq \cO_E$ (elliptic curve is a complex Lie group). The
receptacle is \emph{one-dimensional}, not zero; the non-vanishing
class is the BCOV tree-level Yukawa coupling $Y_3$, which is generically
non-zero on $K3 \times E$. Formality on $K3 \times E$ therefore
proceeds via C\u{a}ld\u{a}raru--Huybrechts $2010$
(arXiv:$0907.2450$ \S$4$) twisted HKR, absorbing Atiyah-class cups into
$\sqrt{\mathrm{td}(TX)}$ under $c_1(X) = 0$; it does \emph{not}
proceed via rank-vanishing of the cubic receptacle.

\emph{Retraction of earlier receptacle identification.} The first-wave
spine stated "Kuranishi cubic lives in
$H^3(K3, \Lambda^3 T_{K3}) = 0$ by rank ($T_{K3}$ rank $2$)". That
vanishing is true on the $K3$ factor alone but is the \emph{wrong
receptacle} for formality on the product $K3 \times E$. The correct
K3 $\times$ E receptacle is one-dimensional and contains the BCOV
Yukawa $Y_3 \neq 0$ generically.

\emph{Retraction of attribution.} The first-wave attribution of the
degree-count formality to "Kapranov $1999$ \S$3.1$" is incorrect:
Kapranov $1999$ \S$3.1$ is the Dolbeault resolution of $\cO_X$, and
\S$4$ Thm.~$2.8.1$ treats the compact-non-flat $L_\infty$-structure
with Atiyah-sourced brackets. The correct primary source for the
flat-space harmonic-retraction formality is Kontsevich $1997$
\S$6.4.1$ lifted to Dolbeault by Costello--Gwilliam $2017$ Vol.~II
Ch.~$7$.
\end{theorem}

\subsection{Miki $S_3$ split across four primary sources}
\label{wn:subsec:second-pass-miki-split}

\begin{theorem}[Four-paper attribution of the $S_3$-triality]
\label{wn:thm:second-pass-miki-four}\ClaimStatusTheorem

The $S_3$-symmetry of the parameter space
$(\epsilon_1, \epsilon_2, \epsilon_3)$ and its descent to the
three-way triality on $Y^+_{\epsilon_1, \epsilon_2, \epsilon_3}
(\widehat{\fgl}_1)$, $\mathrm{CoHA}(\mathbb{C}^3)$, and
$\mathcal{W}_{1+\infty}[\lambda]$ is realised by a chain of four
primary-source theorems:
\begin{enumerate}
\item Miki $2007$ \emph{JMP} $48$:$123520$ Thm.~$1$: proves
$\mathrm{PSL}_2(\ZZ)$-covariance (horizontal $\leftrightarrow$ vertical
exchange on the quantum toroidal $\ddot{U}_{q, \kappa}
(\widehat{\widehat{\fgl}}_1)$); this is \emph{two-axis}
pair-exchange, not full triality.
\item Feigin--Hashizume--Hoshino--Shiraishi--Yanagida $2009$
(arXiv:$0904.1679$) Thm.~$4.3$: establishes the explicit shuffle-kernel
$S_3$-symmetry
$\omega(z, w) = \prod_i (z - w - \epsilon_i)/(z - w)^3$.
\item Feigin--Jimbo--Miwa--Mukhin $2016$
(arXiv:$1603.02765$) Thm.~$2.2$: lifts the shuffle-kernel $S_3$ to a
Hopf-algebra triality on the three quantum-affine subalgebras.
\item Tsymbaliuk $2017$ (arXiv:$1404.5240$): realises the affine
Yangian avatar $Y^+_{\epsilon_1, \epsilon_2, \epsilon_3}
(\widehat{\fgl}_1)$ and the triality descent via the
Schiffmann--Vasserot--Tsymbaliuk three-presentation isomorphism.
\end{enumerate}

\emph{Retraction of single-attribution.} The first-wave spine
attributed the entire triality to "Miki $2007$ $S_3$-automorphism of
$\mathcal{W}_{1+\infty}[\lambda]$". Miki's theorem proves
$\mathrm{PSL}_2(\ZZ)$, not $S_3$; the full triality requires the
four-paper chain above.

\emph{Structural point.} The ambient $S_3$-action on
$\mathrm{Spec}\,\CC[\epsilon_1, \epsilon_2, \epsilon_3]$ is manifest;
every hyperplane $\sigma_1 = \delta$ is $S_3$-stable. The CY slice
$\sigma_1 = 0$ is distinguished only by the central-element constraint
$q_1 q_2 q_3 = 1$ (Miki $2007$), which defines the family of affine
Yangians; the $S_3$-symmetry does not select a slice.
\end{theorem}

\subsection{Three faces of $8$: four corrections to primary-source
attribution}
\label{wn:subsec:second-pass-three-faces}

\begin{theorem}[The three faces of $8$ at the $\mathsf{B}$-row, as a
unified conjecture]
\label{wn:thm:second-pass-three-faces}\ClaimStatusConjectured

The three numerical identifications of $K^{\kch}_{\mathsf{B}} = 8$ at
the Mukai-enhanced K3 Heisenberg $\cH_{\Mukr}(K3)$ on the five-archetype
landmark ceiling are:
\begin{enumerate}[label=(\roman*)]
\item $8 = 2\, c_+(\Mukr(K3)) = 2 \cdot 4$: Mukai signature split of
the lattice $\Wttl(K3) \simeq E_8(-1)^{\oplus 2} \oplus U^{\oplus 3}$
of signature $(4, 20)$. Primary: Mukai $1987$ \emph{Nagoya Math.~J.}
$81$ \S$1$ + Vol~I anomaly-ratio bridge (chiral\_center\_theorem.tex
Theorem~C(a)).
\item $8 = \mathrm{lcm}(2_{\mathrm{mult}}, 4_{\mathrm{Bruinier}})
\cdot 2_{\mathrm{super}}$: monodromy order of
$\cL^{\Delta_5}$ on the $\mu_8$-gerbe banded cover
$(\Delta_5/\eta^{12})^{1/8}$. Primary (decomposed across four papers):
Borcherds $1998$ \emph{J.~reine angew.~Math.} $494$ \S$10$
(paramodular multiplier); Bruinier $2002$ LNM $1780$ Thm.~$5.12$
(divisor formula); Kudla--Millson $1986$ (Arakelov Chern class);
Schauenburg $1998$ (super-parity extension).
\item $8 = \ell$: Lusztig root-of-unity
order $\zeta^{\ell} = 1$ at which the Hall--Drinfeld double specialises
to the small quantum group $\mathbf{u}_{\zeta}$ banded by a
$\mu_8$-gerbe (Lusztig $1990$ \emph{Geom.~Dedicata} $35$;
Kapranov--Voevodsky $1994$; Drinfeld $1990$ quasi-Hopf scaling
$\hbar = 2\pi i / \ell$, with $\hbar^2 K = -1$ valid in the
Kontsevich-torsor normalisation $(2\pi)^2 \mapsto 1$).
\end{enumerate}

\emph{Unpacking the universal trace identity} $\hbar^2 K^{\kappa_{\mathrm{ch}}}
= -1$. The identity reads in three steps: (a) $K^{\kappa_{\mathrm{ch}}}
= \kch(A) + \kch(A^!)$ is the total chiral central charge of the
Koszul pair $(A, A^!)$; for the $\mathsf{B}$-row witness $A =
\cH_{\Mukr}(K3)$ this total equals $8 = 2c_+(\Mukr(K3))$, where $c_+$
is the number of positive-norm directions of the Mukai lattice
(signature $(4, 20)$, so $c_+ = 4$, giving $2c_+ = 8$). The factor of
$2$ arises because both the algebra and its Koszul dual contribute one
copy of $c_+$ through the signature-split trace on the Mukai lattice.
(b) Drinfeld scaling sends the quasi-Hopf associator parameter to
$\hbar = 2\pi i / \ell$; squaring, $\hbar^2 = -(2\pi)^2/\ell^2$. (c)
The Kontsevich-torsor normalisation absorbs the $(2\pi)^2$ factor
into the chosen associator basepoint (the torsor of Drinfeld
associators carries a preferred Kontsevich integral representative
at which $(2\pi)^2$ is absorbed into the trivialisation), leaving
$\hbar^2 = -1/\ell^2$. Substituting $\ell^2 = K^{\kappa_{\mathrm{ch}}}$
(the root-of-unity order squared matches the Koszul total) gives
$\hbar^2 K^{\kappa_{\mathrm{ch}}} = -1$ --- a universal identity
independent of the specific CY family, once the Kontsevich
normalisation is fixed.

The $K = 2c_+$ identification is the geometric heart: $c_+$ counts
positive directions in the Mukai lattice, which is the signature
invariant detecting how much of the Serre functor is self-dual versus
dual-acting. The doubling $K = 2c_+$ reflects that both $\kch(A)$ and
$\kch(A^!)$ each pick up $c_+$ from the respective half of the
Serre bifunctor decomposition; equivalently, $K$ measures the
rank of the Serre kernel restricted to positive directions, which
equals the rank of the mirror-partner sector and produces the
doubling.

Faces (i) $\leftrightarrow$ (ii) are unified by a conjectural
Bruinier--Mukai reciprocity at signature $(4, 20)$
\texttt{conj:BZ-Muk-Br}, verified numerically at three witnesses
(Monster $K = 2$; K3 $K = 8$; Fake Monster $K = 50$). Faces
(ii) $\leftrightarrow$ (iii) are unified by Kapranov--Voevodsky
categorical $\mu_\ell$-gerbe structure: the paramodular multiplier
class of order $\ell$ is precisely the band of the $\mu_\ell$-gerbe
classifying Lusztig's small quantum group, so arithmetic torsion
order equals quantum-group root-of-unity order equals $8$. The
three-faces identity is a \emph{conjectural unification}, not three
independent theorems with automatic reciprocity.

\emph{Retraction of primary-source claims.} The first-wave spine
stated four items that do not hold to literal primary-source
verification:
(a) "Beilinson--Drinfeld Koszul-conductor identity on $(4, 20)$" is
not in BD $2004$ \emph{Chiral Algebras}; correct source is Mukai $1987$
plus Vol~I bridge.
(b) "Bruinier $2002$ Prop.~$5.1$: monodromy $= 2 \cdot 4 = 8$" is a
misquote; Bruinier Prop.~$5.1$ is a local product expansion, not a
torsion-order formula.
(c) "$\eta^9\vartheta_1$ weight $9/2$ contributes index $4$" is a
typographic error; $\eta^9 \vartheta_1$ has weight $5$ index $1/2$.
(d) "Universal $\hbar^2 K = -1$" holds only in Kontsevich-torsor units;
physical Drinfeld scaling gives $-(2\pi)^2 / \ell$.

The ghost-theorem is the conjectural unification above; individual
faces hold unconditionally, their reciprocity is programme-specific.
\end{theorem}

\subsection{Shioda--Inose Mordell--Weil: rank 2 or 16, per fibration
configuration}
\label{wn:subsec:second-pass-shioda-inose-MW}

\begin{theorem}[Shioda--Inose K3 Mordell--Weil by Kodaira
configuration]
\label{wn:thm:second-pass-shioda-inose-MW}\ClaimStatusTheorem

Let $S$ be a K3 surface of Picard rank $\rho(S) = 20$ (Shioda--Inose)
with an elliptic fibration $\pi\colon S \to \bP^1$. By the
Shioda--Tate formula,
\[
\rho(S) \;=\; 2 \;+\; r \;+\; \mathrm{rk}\,\mathrm{MW}(\pi),
\qquad r \;=\; \sum_{v}\,(m_v - 1),
\]
with $r$ the sum of reducible-fibre corrections. Two extremal
configurations:
\begin{itemize}
\item \emph{Two $II^{\ast}$ fibres.} The fibre root lattice is
$L^{\mathrm{(red)}}_{\mathrm{fibre}} = E_8 \oplus E_8$ of rank $16$; by
Shioda--Tate, $\mathrm{rk}\,\mathrm{MW} = 20 - 2 - 16 = 2$. The
Mordell--Weil group has rank $\mathbf{2}$, not $16$.
\item \emph{Two $I_2$ fibres plus $20$ $I_1$ fibres.} The reducible
correction is $r = 2$; $\mathrm{rk}\,\mathrm{MW} = 20 - 2 - 2 = 16$,
and $\mathrm{MW}(\pi) \simeq E_8(-1)^{\oplus 2}$ under a Shioda-height
rescaling by $\chi(\cO_S) = 2$ and subject to conditions on the
transcendental lattice.
\end{itemize}

\emph{Retraction of the combined statement.} The first-wave spine
stated "Shioda--Inose K3 has $\mathrm{MW}(\pi) = E_8 \oplus E_8$" as a
uniform claim. The claim conflates two distinct rank-$16$ lattices:
the fibre root lattice (in the $II^\ast + II^\ast$ configuration,
where $\mathrm{MW}$ is rank $2$) and the Mordell--Weil lattice itself
(in the $I_2 + I_2 + 20\,I_1$ configuration, where
$\mathrm{MW} \simeq E_8(-1)^{\oplus 2}$ with rescaling). Not every
Shioda--Inose K3 carries either fibration; singular K3 surfaces are
parametrised by binary quadratic forms with countably many classes.
The correct claim must specify the Kodaira configuration.

\emph{Ghost-theorem for the elliptic-surface programme.} The
Mordell--Weil lattice, when $\mathrm{MW}(\pi) \simeq E_8(-1)^{\oplus
2}$ at the $I_2 + I_2 + 20\,I_1$ configuration, indexes the
\emph{internal} simple-root system of the Stage-$2$ GBKM
$\fg_{\cE, \bP^1}$ (not an external orbit symmetry), via the Shioda
canonical height plus the rescaling $\chi(\cO_S) = 2$ to match the
$E_8$ Gram matrix. The discriminant-group argument (using
unimodularity of $E_8(-1)^{\oplus 2}$ as a primitive sublattice of the
unimodular $\mathrm{II}_{1,1}\oplus E_8^{\oplus 3}$) makes the
internal-decoration status explicit.
\end{theorem}

\subsection{Fake Monster rank obstruction: tight form at
positive-definite rank}
\label{wn:subsec:second-pass-FM-rank}

\begin{theorem}[Tight positive-rank obstruction for the Fake Monster]
\label{wn:thm:second-pass-FM-rank}\ClaimStatusTheorem

The Fake Monster Lie algebra is excluded from the $d = 3$ Stage-$2$
specialisation by an obstruction on \emph{positive-definite rank} of
the Mukai-lattice-plus-$U(E)$ ambient, not on total rank alone:
\begin{itemize}
\item $\Wttl(K3) = \mathrm{II}_{4, 20}$ has positive rank $4$;
$\Wttl(K3) \oplus U(E) = \mathrm{II}_{5, 21}$ has positive rank $5$.
\item The Leech lattice $\LLeech$ (rank $24$, positive-definite, even
unimodular) requires an ambient with positive rank $\geq 25$ for a
primitive embedding of the enhanced $\mathrm{II}_{25, 1}$ Fake Monster
root lattice.
\item Tight obstruction: $5 < 25$, giving a positive-rank deficit of
$20$; not the loose form $20 < 24$ used in the first-wave spine.
\end{itemize}

At $d = 5$ on $K3 \times K3 \times E$, the ambient
$(\Wttl(K3))^{\otimes 2} \oplus U(E)$ has signature $(417, 161)$ and
positive rank $417$, leaving a positive-rank surplus of $392$ beyond
the Fake Monster requirement. Nikulin $1979$ \emph{Izv.~Akad.~Nauk
SSSR} Thm.~$1.12.2$ guarantees a primitive embedding of $\mathrm{II}_{25,
1}$ into a signature-$(p, q)$ ambient when $p \geq 25$ and $q \geq 1$,
which is satisfied with margin.

Borcherds $1998$ Thm.~$13.3$ applied to
$L = \mathrm{II}_{26, 2}$ with input Jacobi form of weight
$1 - 26/2 = -12$ (matching the weight of $1/\eta^{24}$) yields the
Fake Monster denominator $\Phi_{\mathrm{FM}}$ of Borcherds weight
$c(0)/2 = 24/2 = 12$, where $c(0) = p_{24}(1) = 24$ is the constant
coefficient of $1/\eta^{24}(\tau) = q^{-1}\prod_{n \geq 1}
(1 - q^n)^{-24} = q^{-1}(1 + 24 q + O(q^2))$.

The Borcherds $\Phi_{12}$ on $\mathrm{II}_{26, 2}$ and the Igusa
$\Phi_{12}$ on $\mathbb{H}_2$ are distinct automorphic forms on
distinct domains; the weight coincidence at $12$ is arithmetic, not
structural.
\end{theorem}

\subsection{Class-$\mathcal{S}$ $(A_1, \Sigma_{0, 24})$: upgraded to
theorem at character level}
\label{wn:subsec:second-pass-class-S}

\begin{theorem}[Character-level $c_{4d}(A_1, \Sigma_{0,24}) = 107/6$]
\label{wn:thm:second-pass-class-S-107-6}\ClaimStatusTheorem

The class-$\mathcal{S}$ four-dimensional central charge
$c_{4d}(A_1, \Sigma_{0, 24}) = 107/6$ is a theorem at character level,
via direct pants-decomposition:
\[
c_{4d}(A_1, \Sigma_{0, n}) \;=\; (2\,n_v + n_h)/12
\;=\; (5n - 13)/6.
\]
For $n = 24$: $(n_v, n_h) = (63, 88)$, giving $c_{4d} = 107/6$.
Cross-check: at $n = 4$, $(n_v, n_h) = (3, 8)$, recovering
$c_{4d}(A_1, \Sigma_{0, 4}) = c_{4d}(\mathrm{SU}(2)\text{ with }
N_f = 4) = 7/6$, the Argyres--Douglas value.

\emph{Derivation at pants level.} Decompose $\Sigma_{0, n}$ into
$n - 2$ trinions and $n - 3$ internal $\mathrm{SU}(2)$-tubes. Each
$T_2$ trinion carries $(n_v, n_h)_{T_2} = (0, 4)$ (a half-hyper in the
pseudoreal $(\mathbf{2}, \mathbf{2}, \mathbf{2})$ representation of
$\mathrm{SU}(2)^3$); each internal tube gauges
$\mathrm{SU}(2)$ with central charge $(3, 0)$. Total:
$(n_v, n_h) = (3(n-3), 4(n-2))$, giving the formula above.

\emph{Upgrade from conjecture.} The first-wave spine labelled this
$\Sigma_{0, 24}$ route as conjectural "pending primary-source trace of
the nilpotent-Higgsing combinatorics". The direct pants-level
derivation (using only Gaiotto class-$\mathcal{S}$ + Shapere--Tachikawa
trace anomaly) upgrades the identity to theorem status at character
level. Chacaltana--Distler $2010$ \emph{JHEP} $10$:$099$ Table~$3$
row~$1$ provides the $(5n - 13)/6$ formula; the $n = 4$ cross-check
against Argyres--Douglas $\mathrm{SU}(2)\,N_f = 4$ is independent
verification.

\emph{Uniqueness caveat.} The Diophantine equation
$13(g - 1) + 5n = 107$ has multiple solutions: $(g, n) = (0, 24)$ and
$(5, 11)$ both yield $c_{4d} = 107/6$. The selection of
$\Sigma_{0, 24}$ requires the additional Mukai-lattice / $M_{24}$
constraint $n = 24 = \dim H^*(K3, \ZZ)$, which remains conjectural at
the CY$_3$ level.
\end{theorem}

\subsection{Prochazka--Rapcak: surjection not embedding, and the corner
Yangian}
\label{wn:subsec:second-pass-PR}

\begin{theorem}[Prochazka--Rapcak truncation, with the corner-Yangian
translation]
\label{wn:thm:second-pass-PR}\ClaimStatusTheorem

Prochazka--Rapcak $2018$ (arXiv:$1711.09993$) Thm.~$1.1$ constructs a
\emph{surjection}
\[
\cW_{\infty}[\mu, \lambda] \;\twoheadrightarrow\; \cW_N(\fsl_N)_\kappa,
\qquad\text{at }\mu = N,
\]
from the universal two-parameter $\cW_{\infty}$ to the principal
$\cW_N$ at specialised $\mu$. This is a \emph{quotient}, not an
embedding.

\emph{Yangian avatar.} Under the Schiffmann--Vasserot--Tsymbaliuk
three-presentation isomorphism, the PR truncation lifts to a Yangian
surjection
\[
Y_{\epsilon_1, \epsilon_2, \epsilon_3}(\widehat{\fgl}_1)
\;\twoheadrightarrow\; \mathrm{Heis} \;\otimes\;
Y_{\epsilon_1, \epsilon_2}(\widehat{\fsl}_N)_{\kappa(N)},
\]
a single $\fsl_N$-Yangian tensored with Heisenberg — \emph{not} a
tensor power $Y(\widehat{\fgl}_1)^{\otimes n}$ with
level-$n$ identification.

\emph{Retraction.} The first-wave spine stated the direction
backwards, as an embedding $Y(\widehat{\fsl}_n) \hookrightarrow
Y(\widehat{\fgl}_1)$; the direction is surjection, not embedding. A
second latent error in the earlier inscription at
\texttt{chapters/theory/quantum\_chiral\_algebras.tex}
lines~$4025$--$4053$ uses the tensor-power form
$Y_{0, 0, n}(\widehat{\fsl}_n) \hookrightarrow
Y_{0, 0, 1}(\widehat{\fgl}_1)^{\otimes n} /
(\text{level-}n\text{ identification})$,
attributed to "Prochazka--Rapcak $2018$ \S$5$"; this is likewise wrong
(the tensor target is a physically different configuration of
$n$ non-interacting $\mathrm{GL}(1)$-Wilson lines versus a single
$\mathrm{GL}(n)$-Wilson-line corner). The correct corner-Yangian
translation is the single $\fsl_N$-Yangian + Heis form above.

\emph{Chiral-Yangian scope.} The earlier "chiral Yangian
$\cY^{\mathrm{ch}}(\fg)$ is the $L_0$-bounded sub-VOA of the corner"
phrasing is internally incoherent: principal $\cW$-algebras at generic
level have no proper $L_0$-bounded sub-VOAs. The correct three-way
distinction is:
(affine Yangian $Y(\widehat{\fg})$) $\neq$ (principal $\cW$-algebra
$\cW_N(\fg)_\kappa$) $\neq$ (chiral Yangian $\cY^{\mathrm{ch}}(\fg)$);
the relation among them is the PR surjection (or its Yangian lift),
not sub-VOA inclusion.
\end{theorem}

\subsection{AdS$_3$ Brown--Henneaux at the CHL orbifold}
\label{wn:subsec:second-pass-AdS3-BH}

\begin{theorem}[$c_L^{\mathrm{reduced}} = 24$ uniformly via volume
absorption]
\label{wn:thm:second-pass-AdS3-BH}\ClaimStatusTheorem

The reduced Brown--Henneaux central charge of the chiral boundary of
the IIB $D_1$-$D_5$-KK near-horizon geometry
$\mathrm{AdS}_3 \times S^3/\ZZ_N \times K3^{g_N}$ is
\[
c_L^{\mathrm{reduced}} \;=\; 24
\]
uniformly in $N$. The mechanism is a volume-absorption into the
effective three-dimensional Newton constant: Brown--Henneaux $1986$
gives $c = 3\ell / (2G_N^{(3)})$ as an asymptotic-symmetry invariant,
and the $K3^{g_N}$ orbifold volume-factor $24/(N+1)$ is absorbed into
$G_N^{(3)}$ by MSW dimensional-reduction gauge fixing.

\emph{The volume absorption in four steps.} (a) The ten-dimensional
IIB Newton constant $G_N^{(10)}$ dimensionally reduces to three
dimensions by integrating over the $S^3 \times K3^{g_N}$ internal
geometry: $1/G_N^{(3)} = \mathrm{Vol}(S^3) \cdot \mathrm{Vol}(K3^{g_N})
/ G_N^{(10)}$. (b) The CHL orbifold scales the K3-volume by
$24/(N+1)$ (this is the Euler characteristic of $K3^{g_N}$ as a
$\ZZ_N$-orbifold, equal to $\chi(K3)/N + \text{fixed-point correction}
= 24/(N+1)$ after the McKay-Thompson count). (c) Substituting into
Brown--Henneaux gives $c_L(N) = 3\ell \cdot 24/(N+1) / (2 G_N^{(3,0)})$
where $G_N^{(3,0)}$ is the $N = 0$ reference Newton constant. (d)
The MSW dimensional-reduction gauge fixing chooses
$G_N^{(3)}(N) = G_N^{(3,0)} \cdot 24/(N+1)$ so the factor absorbs,
giving $c_L^{\mathrm{reduced}} = 24$ uniformly. The point: the
$N$-dependence cancels between the CHL orbifold volume and the MSW
gauge fixing, leaving a universal $24$.

\emph{Two CHL sets, two paramodular-weight scopes.} The parameter
$k_N$ appearing in the Sen dyon-count formula is the Borcherds weight
of the paramodular lift $\Phi_{k_N}$, not a central charge:
\begin{itemize}
\item Gritsenko--Nikulin CHL set $\{1, 2, 3, 4, 6\}$ (those $N$ with
$\varphi(N) \mid 2$ admitting a paramodular witness): weights
$(k_N) = (5, 2, 1, 1, 1)$ (cf.~the single-ladder theorem
\ref{wn:thm:second-pass-single-ladder}).
\item Jatkar--Sen CHL set $\{1, 2, 3, 5, 7\}$ (David--Jatkar--Sen
$2006$ \emph{JHEP} $0606$:$064$): weights $k(N) = 24/(N+1) - 2 \in
\{10, 6, 4, 2, 1\}$.
\end{itemize}
The two sets are different CHL classifications and must be named at
every use. The MSW relation $c_L = 6 k + 24$ uses the Jatkar--Sen
weight as the M$5$-stack level $k$, \emph{not} the Gritsenko--Nikulin
weight.

\emph{Graviton-finiteness statement scope.} The identification
``graviton finiteness $\equiv E_2$-chiral rigidity $\equiv
\mathrm{HH}^2_{E_2}$-vanishing'' is a theorem at $N = 1$
(Borcherds $1986$ Euler-product denominator +
Maloney--Witten $2010$ pure-gravity partition function); at
$N \geq 2$ it requires a $g_N$-equivariant no-ghost refinement
currently not in primary literature and is
\ClaimStatusConjectured.
\end{theorem}

\subsection{Two-module retraction at the affine Yangian of $\fgl_1$}
\label{wn:subsec:second-pass-two-modules}

\begin{theorem}[$\mathrm{Hilb}(\mathbb{C}^2)$-Fock and
$\cM_n(\mathbb{C}^3)$-CoHA are inequivalent modules]
\label{wn:thm:second-pass-two-modules}\ClaimStatusTheorem

The affine Yangian of $\fgl_1$ carries two distinct modules, with
different characters:
\begin{itemize}
\item $\cF_{\mathrm{Nak}} \;=\; \bigoplus_n H^{\ast}_T(\mathrm{Hilb}^n
(\mathbb{C}^2))$: the rank-$1$ Nakajima Fock module on the smooth
Hilbert scheme of points on $\mathbb{C}^2$, character
$\chi(\cF_{\mathrm{Nak}})(q) = \prod_{n \geq 1}(1 - q^n)^{-1}$ (the
partition-function generating function $1/(q; q)_\infty$). The three
equivariant parameters enter via the shuffle kernel with
$\epsilon_3 = -\epsilon_1 - \epsilon_2$ attached to the normal
direction.
\item $\cM_{\mathrm{reg}} \;=\; \bigoplus_n H^{\ast}_T(\cM_n
(\mathbb{C}^3), \phi_W)$: the CoHA regular representation on the
Donaldson--Thomas vanishing-cycle cohomology of the Jordan-triple
quiver moduli, character $\chi(\cM_{\mathrm{reg}})(q) = M(-q) =
\prod_{n \geq 1}(1 - (-q)^n)^{-n}$ (the MacMahon generating function).
\end{itemize}

These modules are \emph{not equivalent}: fixed-point counts differ
from $n = 2$ (partitions $p(2) = 2$ versus plane partitions
$p_3(2) = 3$). They are structurally distinct representations in
different roles.

\emph{Correction to first-wave spine.} The first-wave retraction \#$14$
stated "the correct module is $\bigoplus_n H^{\ast}_T(\mathrm{Hilb}^n
(\mathbb{C}^2))$, equivalently $\bigoplus_n H^{\ast}_T(\cM_n
(\mathbb{C}^3), \phi_W)$". The two are \emph{not} equivalent. The
direction of the retraction is correct (ruling out
$\mathrm{Hilb}^n(\mathbb{C}^3)$ via Iarrobino--Brian\c con
singularity), but both modules play distinct roles in the affine-Yangian
module theory.

\emph{Typographic correction.} At
\texttt{working\_notes.tex} lines $1047$--$1091$, the
"$K_T(\mathrm{Hilb}^n(\mathbb{C}^3))$" module label should read
"$H^{\ast}_T(\mathrm{Hilb}^n(\mathbb{C}^2))$" (the
$R$-matrix formula uses the $2$-d Young-diagram content
arm/leg/co-arm, which is on $\mathrm{Hilb}^n(\mathbb{C}^2)$);
the surrounding mathematics is otherwise correct.
\end{theorem}

\subsection{Wave-$1$ promotions: three theorems from the adversarial
residue}
\label{wn:subsec:second-pass-promotions}

\begin{theorem}[Three Wave-$1$ frontier items promoted to theorem]
\label{wn:thm:second-pass-promotions}\ClaimStatusTheorem

The following first-wave frontier items are upgraded to theorem
status on the basis of the second-wave primary-source audit:
\begin{enumerate}
\item \emph{BCFG anomaly universal vanishing}:
$\kanom(X, \fg^{\mathrm{BCFG}}) = 0$ for every CY$_3$ $X$ and every
non-simply-laced $\fg \in \{B_n, C_n, F_4, G_2\}$, by a
uniform folding-based proof: $\sigma^{\ast} = -1$ on the
one-dimensional $S^3(\fsl_n)^{\fsl_n}$-invariant space, so $\sigma$-fixed
points of $d^{abc}$ vanish. Quadratic Casimirs explicit:
$C_2(B_n) = 2n - 1$, $C_2(C_n) = n + 1$, $C_2(F_4) = 9$,
$C_2(G_2) = 4$. Promotes earlier conjectural BCFG extension of the
ADE all-orders theorem.
\item \emph{Class-$\mathcal{S}$ $c_{4d}(A_1, \Sigma_{0, 24}) = 107/6$}:
theorem at character level via direct pants-decomposition
(Theorem \ref{wn:thm:second-pass-class-S-107-6}). Cross-checked at
$n = 4$ against the Argyres--Douglas $\mathrm{SU}(2)\,N_f = 4$ value
$7/6$.
\item \emph{$\mathsf{B}$-row $K^{\kch} = 8$ via direct Serre bifunctor}:
unconditional on the Bruinier Heegner-Chern-class black box, using
only the Serre functor $S_{K3} = [2]$ acting on the derived bar
coalgebra of $\cH_{\Mukr}(K3)$, giving $\mathrm{ord}(S_{K3}^2 \otimes
\tau_E) = 8$ in $\mathrm{Heisdouble}(K3 \times E)$. The three faces
(Mukai / Humbert / Lusztig) are then three functorial projections of
this single structural identification, rather than three independent
identifications unified post-hoc by Bruinier reciprocity.
\end{enumerate}
\end{theorem}

\subsection{Three-tier frontier stratification}
\label{wn:subsec:second-pass-frontier-tiers}

\begin{theorem}[Three tiers of residual frontier]
\label{wn:thm:second-pass-frontier-tiers}

The residual frontier stratifies into three tiers by the character of
the primary-source gap closing the item:

\emph{Tier I — tight (single named primary-source gap).}
\begin{itemize}
\item BCFG $\sigma$-equivariant renormalisation scheme for Costello
$6$d $\hCS$ (Costello--Gwilliam Vol.~II \S$11.1$ gap).
\item Elliptic-surface specialisation
$(\Sigma_2, C) = (\mathcal{E}, \bP^1)$ at Shioda--Inose $\rho = 20$
(Borcherds $1998$ Thm.~$13.3$ on signature $(2, 20)$).
\item Ran-level Miki $S_3$-triality as factorisation-algebra
automorphism (Costello--Paquette $2020$ \S$5$).
\item Bardeen--Zumino cochain bridging consistent and covariant
cubic-Casimir anomaly representatives as an $L_\infty$-morphism
($N_1$).
\end{itemize}

\emph{Tier II — moderate (method extension).}
\begin{itemize}
\item 3-dualizability on compact CY$_3$ (gated by Tier III integral
$E_d$-formality).
\item Bracket-level $\fg_{\mathrm{BPS}}(K3 \times E) \simeq
\fg_{\Delta_5}$, reducing to a single Hecke--Borcherds identity on the
Gritsenko $1999$ paramodular family.
\item Dimension-stratified GKM census $\{\fg^{\mathrm{BKM}}_d\}$
across $d \in \{3, 5\}$ unifying Monster (virtual CY$_3$),
Fake Monster ($d = 5$, $\kBKM = 12$), and $\fg_{\Delta_5}$ ($d = 3$,
$\kBKM = 5$) via Borcherds-functorial restriction + Nikulin
primitive-embedding input ($N_2$).
\item Three-faces-of-$8$ unification at
$\mathrm{Heisdouble}(K3 \times E)$ via the single structural
identification $\mathrm{ord}(S_{K3}^2 \otimes \tau_E) = 8$ ($N_3$).
\end{itemize}

\emph{Tier III — loose (new machinery required).}
\begin{itemize}
\item Integral $E_d$-formality at $d \geq 3$.
\item $(\infty, 1)$-functoriality of $\Phi^{\mathrm{FA}}_d$ on
non-formal CY categories.
\item Non-abelian Fake Monster doubly-reduced Donaldson--Thomas
integrand $= 1/\Phi_{12}$ on $K3 \times K3 \times E$.
\item Non-CHL $N = 7$ order-$4$ central extension of $\Mpr_4$ by
$\mu_4$.
\item Rank-$\geq 3$ lattice-polarised $\fg_L$ family.
\end{itemize}
\end{theorem}

\subsection{The picture in one sentence, refined}
\label{wn:subsec:second-pass-one-line}

\emph{A single $E_d^{\mathrm{hol}}$-factorisation algebra
$\cF_X$ on each Calabi--Yau $d$-fold $X$ is the canonical Stage-$1$
output of $\Phi$, realised at $d = 3$ as the Costello--Francis--Gwilliam
$E_3^{\mathrm{hol}}$-algebra of observables of six-dimensional
holomorphic Chern--Simons on $X$; transverse-bordism classes of
$(\Sigma_{d-1}, C)$-pairs specialise $\cF_X$ to $E_1$-chiral shadows,
of which Igusa $\Delta_5$ GBKM $\fg_{\Delta_5}$ at $d = 3$, the virtual
Borcherds Monster on $\mathrm{II}_{1, 1}$, and the Fake Monster at
$d = 5$ on $K3 \times K3 \times E$ with $\kBKM = 12$ are three
representatives of a dimension-stratified Borcherds--Kac--Moody census
unified by the universal identity $\kBKM(\Phi) = c(0)/2$ and by
Wang--Williams $2023$ pullback rigidity of $\Phi_{12}$; the CHL
Borcherds-weight ladder $(5, 2, 1, 1, 1)$ at $N \in \{1, 2, 3, 4, 6\}$
is the output of three mutually-compatible primary lifts rather than
of two distinct ladders; the boundary rows $N \in \{5, 7, 8\}$ arise
as metaplectic paramodular forms on $\Mpr_4(\mathbb{A})$ via
Niwa $1974$ combined with the Howe pair $(\Mpr_4, \mathrm{O}(1, 1))$;
on $K3 \times E$ specifically the six routes to the quantum vertex
chiral group $G(K3 \times E)$ are three $(\Sigma_2, C)$-specialisations
of $\cF_{K3 \times E}$ plus three routes consuming non-cycle-class data,
collectively indexed by the triple (orbit, construction machine,
generator-rank), with generator-rank $\rho^{R_i} \in \{3, 12, 24\}$;
the four-value crystallisation $\{2, 3, 5, 24\}$ is the monograph
convention naming six invariants across two Calabi--Yau categories
via a four-functor square, with $\mathsf{B}$-row witness
$K^{\kch}_{\mathsf{B}} = 8$ as the single structural identification
whose three faces (Mukai, Humbert, Lusztig) are functorial projections
of $\mathrm{ord}(S_{K3}^2 \otimes \tau_E) = 8$ on
$\mathrm{Heisdouble}(K3 \times E)$.}

% End of second-pass refinement.

% =====================================================================
% Third-pass Bezrukavnikov-voice refinement: D-module / Shimura / MGSL
% / Chenevier / Humbert rigour, six attack-heal cycles.  The second
% adversarial pass exposed category slips at the level of invariants
% and attributions; the third pass attacks the representation-theoretic
% substrate -- D-module categories, derived-category base data, Shimura
% reductive-group-and-level data, motivic Galois coefficient rings,
% Chenevier determinant conditions, and Humbert-divisor normalisations.
% Each cycle states the Bezrukavnikov attack, the corrected statement,
% and the primary-source anchors with the hypotheses fully named.
% =====================================================================

\section{Third-pass refinement: geometric representation theory at
Bezrukavnikov precision}
\label{wn:sec:spine-third-pass}

\emph{The second-pass refinement corrected category attributions and
primary-source traces; a third adversarial pass against the
representation-theoretic substrate exposes statements whose categories
of sheaves, reductive-group/level data, Hodge-structure types, and
coefficient rings need to be named at every use for the theorem to be
precise rather than merely true. The attacks below sort into six
cycles: D-module and derived-category categories, Shimura variety and
Hodge-structure data, motivic Galois super-Lie algebra coefficient
rings, Chenevier determinant versus pseudo-character,
Humbert-divisor discriminant conventions, and class-$\mathcal{S}$
gauge-and-flavour content plus the $8\times 5$ bistrata
$(\infty,1)$-categorical framework. Each cycle inscribes a
heal that fixes the category, hypothesis, or base ring at the
precision required for the theorem to stand on primary-source
foundations (Beilinson--Bernstein $1981$; Saito $1990$;
Hotta--Takeuchi--Tanisaki $2008$; Drinfeld--Gaitsgory $2015$;
Chenevier $2014$; van~der~Geer $1988$).}

\subsection{Cycle 1: D-module and derived-category categories, fully
named}
\label{wn:subsec:third-pass-D-module}

\begin{theorem}[Derived categories with base data named]
\label{wn:thm:third-pass-derived-base}\ClaimStatusTheorem

The derived categories invoked in the $K3 \times E$ four-functor
square (Theorem~\ref{wn:thm:second-pass-four-functor-square}) are:
\[
D^{\mathrm{b}}_{\mathrm{coh}}(K3, \mathcal{O}_{K3}),
\qquad
D^{\mathrm{b}}_{\mathrm{coh}}(K3 \times E, \mathcal{O}_{K3\times E}),
\]
the bounded derived categories of complexes of quasi-coherent
$\mathcal{O}$-modules with coherent cohomology, over the structure
sheaf of the corresponding complex projective variety. The base ring
is $\CC$; the category is $k$-linear over $k = \CC$; the bounded
structure is on cohomological amplitude, and coherence requires
finite generation of every cohomology sheaf over $\mathcal{O}$.

On the K3 factor the Serre functor is
$S_{K3}(-) = (-)\otimes_{\mathcal{O}_{K3}} \omega_{K3}[2] = (-)[2]$
(shift by complex dimension since $\omega_{K3} \simeq \mathcal{O}_{K3}$
for CY$_2$). On $K3\times E$ the Serre functor is
$S_{K3\times E}(-) = (-)[3]$. Serre duality on these categories is
Bondal--Kapranov $1989$ (\emph{Math.~USSR Sbornik}~$70$:$93$),
recalled in Huybrechts $2006$ Ch.~$3$; the Calabi--Yau shift is
precisely the complex dimension $d$ of the variety with trivial
canonical bundle.

\emph{Bezrukavnikov attack on the first-wave prose.} The first-wave
spine wrote ``$D^b\,\mathrm{Coh}(K3)$'' without naming the base ring,
without specifying that the coefficient is $\mathcal{O}$-modules, and
without distinguishing the bounded-derived from the derived-category-
of-quasicoherent-sheaves $D(\mathrm{QCoh}(K3))$, which differ on
unbounded objects relevant to Fourier--Mukai-compatible extensions.
\emph{Heal.} Every invocation of the derived category in this
monograph names $(K3, \mathcal{O}_{K3})$ or
$(K3 \times E, \mathcal{O}_{K3\times E})$, $k = \CC$, and specifies
bounded coherent unless otherwise stated. Where the Fourier--Mukai
transform extends to unbounded complexes, the category
$D(\mathrm{QCoh})$ is named explicitly.
\end{theorem}

\begin{theorem}[Six-dimensional holomorphic Chern--Simons
factorisation algebras, D-module category named]
\label{wn:thm:third-pass-hCS-D-module}\ClaimStatusTheorem

The holomorphic factorisation algebra $\cF_X$ attached by Stage~$1$
to a Calabi--Yau $3$-fold $X$ takes values in the category of
left $\mathcal{D}_{X^n}$-modules on $\mathrm{Ran}(X)$ in the
sense of Beilinson--Drinfeld $2004$ \emph{Chiral Algebras} Ch.~$3.3$
and Francis--Gaitsgory $2012$, with the $(\infty,1)$-category refined
to the presentable stable $\infty$-category
$D_{\mathrm{hol}}^{\mathrm{R}}(\mathcal{D}_{X^n})$ of left
$\mathcal{D}$-modules with regular singularities (Bernstein
$1972$; Kashiwara $1984$; Drinfeld--Gaitsgory $2015$
\emph{Geom.~Funct.~Anal.}~$23$:$149$). Over the formal-disk skeleton
the $\mathcal{D}$-module structure restricts to a
jet-$\mathcal{D}_{\infty}$-module, with Beilinson--Drinfeld $2004$
Prop.~$3.3.3$ identifying the factorisation-algebra axioms with
$\mathcal{D}$-module compatibilities for the Ran symmetric-product
diagonal.

\emph{Bezrukavnikov attack.} The first-wave spine stated
``$\cF_X \in E_3^{\mathrm{hol}}\text{-}\mathrm{HolFA}(X)$'' without
naming (a) which $\mathcal{D}$-module category (holonomic, coherent,
regular-holonomic, quasi-coherent), (b) whether left or right
$\mathcal{D}$-module, (c) the base-ring-and-topology data
($k = \CC$, analytic versus algebraic $\mathcal{D}$-module sheaf).
These are not bookkeeping: left vs.\ right toggles the Koszul-duality
direction (Beilinson--Ginzburg--Soergel $1996$
\emph{J.~Amer.~Math.~Soc.}~$9$:$473$), and holonomic vs.\
regular-holonomic toggles the Riemann--Hilbert correspondence
validity (Kashiwara $1984$).

\emph{Heal, two-lane discipline.} In the chain-level lane,
$\cF_X$ is a prefactorisation algebra of chain complexes of
left regular-holonomic $\mathcal{D}_{X^n}$-modules on the Ran
space over $\CC$, satisfying the Costello--Gwilliam locality axioms;
this is Bernstein $1972$ plus Beilinson--Drinfeld $2004$ Ch.~$3$. In
the $(\infty,1)$-categorical lane, $\cF_X$ is an object of the
symmetric-monoidal presentable stable $\infty$-category
$D_{\mathrm{hol}}^{\mathrm{R, left}}(\mathcal{D}_{\mathrm{Ran}(X)})$
of Drinfeld--Gaitsgory $2015$ Thm.~$6.8.1$, with the $E_3^{\mathrm{hol}}
$-operadic structure realised as a commutative-algebra structure in
the chiral-categorical tensor product
$\otimes^{\mathrm{ch}}_{\mathrm{Ran}(X)}$. Regular-holonomicity gates
the Riemann--Hilbert comparison
$D_{\mathrm{rh}}^{\mathrm{left}}(\mathcal{D}_X)
\xrightarrow{\sim} D^{\mathrm{b}}_{\mathrm{c}}(X(\CC)^{\mathrm{an}}, \CC)$
between left regular-holonomic $\mathcal{D}$-modules and constructible
sheaves on the complex-analytic space, cited as Kashiwara $1984$
(\emph{Publ.~RIMS}~$20$:$319$); without regular singularities, this
equivalence fails.
\end{theorem}

\begin{theorem}[Perverse-sheaf category for the Saito Hodge module
pullback]
\label{wn:thm:third-pass-perverse-saito}\ClaimStatusTheorem

The mixed Hodge module pullback invoked implicitly in the
Hodge-supertrace derivation of $\kch = 0$ on compact CY$_d$ requires
the Saito category $\mathrm{MHM}(X)$ of polarisable mixed Hodge
modules on a smooth complex projective variety $X$ of complex
dimension $d$ (Saito $1990$ \emph{Publ.~RIMS}~$26$:$221$). The
forgetful functor
\[
\mathrm{rat}\colon \mathrm{MHM}(X) \longrightarrow
\mathrm{Perv}(X(\CC)^{\mathrm{an}}, \QQ),
\]
lands in the abelian category of perverse sheaves on the
complex-analytic space with rational coefficients, with heart the
$\mathsf{p}^{+}$-perverse $t$-structure of Beilinson--Bernstein--
Deligne $1982$ (\emph{Ast\'erisque}~$100$). The Hodge supertrace
identity for $\kch$ is realised at the level of the associated graded
for the weight filtration,
\[
\kch(X) = \sum_{p, q} (-1)^{p + q}\,
\mathrm{dim}_{\QQ}\, \mathrm{Gr}^W_{p+q}\,
H^{p + q}_{c}(X, \mathrm{rat}(\QQ_X^H))^{p, q},
\]
restricted to the $(p, 0)$-components via the CY$_d$ Hodge-symmetry
$h^{p, 0}(X) = h^{d - p, d}(X)$
and the trivial-canonical condition $\omega_X \simeq \mathcal{O}_X$.

\emph{Bezrukavnikov attack.} The first-wave spine's ``Hodge
supertrace'' language did not name (a) the perverse-sheaf
$t$-structure ($\mathsf{p}^{+}$ versus $\mathsf{p}^{-}$, which toggle
duality and affect signs), (b) the coefficient ring (is this $\QQ$,
$\CC$, or a Hodge-structure with a polarisation over $\ZZ$?),
(c) the weight-filtration grading (mixed vs.\ pure), (d) the
compactness or smoothness hypothesis governing whether the six-functor
formalism on $\mathrm{MHM}$ applies without defects. \emph{Heal.} The
Hodge-supertrace theorem applies on smooth projective CY$_d$ varieties
$X/\CC$ with $\omega_X \simeq \mathcal{O}_X$; the category is
$\mathrm{MHM}(X, \QQ)$ with forgetful to $\mathrm{Perv}(X^{\mathrm{an}},
\QQ)$ using the $\mathsf{p}^{+}$-perversity in Saito's convention;
the coefficient ring for the Hodge structure is $\QQ$; the
compactness hypothesis enters via Saito's relative decomposition
theorem for proper maps (Saito $1990$ Thm.~$0.1$; Saito $1990$b
\emph{Publ.~RIMS}~$26$:$921$ \S$2.14$). Without compactness, the
supertrace identity requires $c$-supports.
\end{theorem}

\subsection{Cycle 2: Shimura variety and Hodge structure data, fully
named}
\label{wn:subsec:third-pass-Shimura}

\begin{theorem}[Kuga--Satake construction on K3 with full Hodge datum]
\label{wn:thm:third-pass-kuga-satake}\ClaimStatusTheorem

Let $(X, \phi)$ be a polarised K3 surface over $\CC$ of degree
$2e = \phi^2 \in 2\ZZ_{\geq 1}$, with transcendental lattice
$T(X) \subset H^2(X, \ZZ)$ of rank $22 - \rho(X)$ and Hodge
structure of weight $2$ and type $(1, 20, 1)$ on the primitive
cohomology $P^2(X, \QQ) = \phi^{\perp} \subset H^2(X, \QQ)$, with
Picard rank $\rho(X) = \mathrm{rk}\,\mathrm{NS}(X)$. The Kuga--Satake
construction is the functor
\[
\mathrm{KS}\colon \bigl(\mathrm{K3}_{(2e)}\bigr)_{\CC}^{\mathrm{pol}}
\longrightarrow \mathrm{Sh}_{K}\bigl(\mathrm{GSpin}\bigl(P^2(X,
\QQ)\bigr)\bigr)_{\CC},
\]
where the source is the moduli of degree-$2e$ polarised complex K3
surfaces and the target is the complex points of the Shimura variety
of level structure $K \subset \mathrm{GSpin}(P^2)(\mathbb{A}_f)$
attached to the reductive group $\mathrm{GSpin}(P^2(X, \QQ))$,
the general spin group of the primitive-cohomology quadratic form
over $\QQ$. Primary: Kuga--Satake $1967$ \emph{Illinois J.~Math.}~$11$;
Deligne $1972$ \emph{Invent.~Math.}~$15$:$206$; Rizov $2010$
\emph{J.~reine~angew.~Math.}~$648$.

\emph{Target Shimura variety type.} The target is the Shimura variety
of Hodge type (Deligne $1979$ \emph{Proc.~Symp.~Pure~Math.}~$33$:$247$)
attached to the pair $(\mathrm{GSpin}(P^2), h)$ with Hodge cocharacter
$h\colon \mathbb{S} \to \mathrm{GSpin}(P^2)_{\RR}$ encoding the weight-$2$
Hodge structure on the primitive-cohomology $P^2(X, \RR)$. It is
not the Siegel Shimura variety $\mathcal{A}_g$ directly; the
embedding into a Siegel datum goes via the Clifford
representation $\mathrm{GSpin}(P^2) \hookrightarrow \mathrm{GSp}(V_{\mathrm{Cl}})$,
where $V_{\mathrm{Cl}} = \mathrm{Cl}^+(P^2)$ is the even Clifford
algebra with its natural symplectic pairing; the image of the K3
inside $\mathcal{A}_g$ is a Shimura subvariety of Hodge type, not
the whole Siegel space.

\emph{Level structure.} The level subgroup
$K \subset \mathrm{GSpin}(P^2)(\mathbb{A}_f)$ is determined by the
discriminant-form structure on $T(X)^\vee/T(X)$ and the polarisation
class; for generic polarisation $e = 1$ on the transcendental lattice
$T(X) = \mathrm{II}_{2, 20}$ the canonical level is the stabiliser of
the genus-form, a finite-index open-compact subgroup of
$\mathrm{GSpin}(\mathbb{A}_f)$. The level must be named at every use;
different levels produce non-isomorphic Shimura varieties.

\emph{Bezrukavnikov attack.} The first-wave spine wrote
``$\mathrm{KS}\colon \mathrm{CY}_2^{\mathrm{Siegel-aut}}
\to \mathrm{AbVar}^{\mathrm{Shimura}}$'' without naming (a) the
reductive group on the source (no ``Siegel-auto'' attached to the
source; the Siegel lives on the target), (b) the target Shimura
variety type (Siegel $\mathcal{A}_g$ versus Hodge-type
$\mathrm{Sh}_K(\mathrm{GSpin})$ -- genuinely different), (c) the level
$K$, (d) the polarisation degree $2e$ that determines the connected
component. \emph{Heal, with Bezrukavnikov precision.} The source is
the moduli of degree-$2e$ polarised K3 surfaces over $\CC$
(a Deligne--Mumford stack); the target is
$\mathrm{Sh}_K(\mathrm{GSpin}(P^2(X, \QQ)))_{\CC}$, a Shimura variety
of Hodge type; the embedding into Siegel $\mathcal{A}_g$ factors
through the Clifford representation; the level $K$ is the
arithmetic-subgroup stabiliser of the discriminant form plus
polarisation; all four data must be named at every invocation of
``Kuga--Satake''.
\end{theorem}

\begin{theorem}[Siegel modular form data with weight and level]
\label{wn:thm:third-pass-siegel-weight-level}\ClaimStatusTheorem

The automorphic forms invoked in the Igusa / Borcherds / Gritsenko
dictionary live on the following group-weight-level triples:
\begin{align*}
\Delta_5 &\in S_5(K(1), \nu_{\Delta_5}),
& K(1) &\subset \mathrm{Sp}_4(\QQ)\text{ paramodular of level }1; \\
\Phi_{10} &\in S_{10}(\mathrm{Sp}_4(\ZZ)),
& \text{level }1\text{ of full Siegel}; \\
\Phi_{12}^{\mathrm{Borch}} &\in M_{12}(\mathrm{O}^+(\mathrm{II}_{26, 2}),
\chi_{12}),
& \text{Fake-Monster domain, trivial level}.
\end{align*}
Every Siegel modular form invocation in the monograph names the
triple $(\text{reductive group}, \text{weight}, \text{level},
\text{multiplier})$. The paramodular group $K(t)$ at level $t$
is the stabiliser of the non-principal polarisation
$\mathrm{diag}(1, t)$ in $\mathrm{Sp}_4(\QQ)$, with $K(1) =
\mathrm{Sp}_4(\ZZ)\cdot \langle V_1 \rangle$ where $V_1$ is the
Fricke involution (Gritsenko $1999$ \S$1.1$; Poor--Yuen $2015$
\emph{Acta Arith.}~$170$:$1$). Without the level $t$ specified,
``paramodular form'' is ambiguous.

\emph{Bezrukavnikov attack.} The first-wave spine frequently wrote
``paramodular $\Delta_5$'' without naming the paramodular level
$t = 1$, and wrote ``$\Phi_{12}$'' without distinguishing the
Borcherds paramodular $\Phi_{12}$ on $\mathrm{O}^+(\mathrm{II}_{26,
2})$ from the Igusa cusp form $\Phi_{12}^{\mathrm{Ig}}$ of weight $12$
on $\mathrm{Sp}_4(\ZZ)$ (Igusa $1962$
\emph{Amer.~J.~Math.}~$84$:$175$). These are two distinct
automorphic forms on two distinct reductive groups with two distinct
weight-twelve coincidences; the ``weight $12$'' collision is
arithmetic, not structural. \emph{Heal.} Wang--Williams $2023$
Thm.~$3.1$'s rigidity statement is for the Borcherds paramodular
$\Phi_{12}^{\mathrm{Borch}}$ on $\mathrm{II}_{26, 2}$; the Igusa
$\Phi_{12}^{\mathrm{Ig}}$ on $\mathrm{Sp}_4(\ZZ)$ is a separate
object. Every invocation of ``$\Phi_{12}$'' names the reductive
group and domain.
\end{theorem}

\begin{theorem}[Shimura--Waldspurger at metaplectic weight with full
level data]
\label{wn:thm:third-pass-metaplectic-level}\ClaimStatusTheorem

The metaplectic paramodular forms at $N \in \{5, 7, 8\}$
(Theorem~\ref{wn:thm:second-pass-boundary-extension}) live on the
genuine (non-split) twofold metaplectic cover
$\widetilde{\mathrm{Sp}_4}(\mathbb{A}_{\QQ})$, with level subgroup
$\widetilde{K}(N) \subset \widetilde{\mathrm{Sp}_4}(\mathbb{A}_f)$
the preimage of the paramodular $K(N)$ under the cover. The Niwa
$1974$ \emph{Nagoya Math.~J.}~$56$:$147$ correspondence
\[
\mathrm{Niwa}\colon S_{k + 1/2}(\Gamma_0(4N), \chi)
\longrightarrow S_{2k - 1}(\Gamma_0(2N), \chi^2)
\]
is a theta-lift realising half-integer to integer weight elliptic
cusp forms, using the Howe dual pair
$(\widetilde{\mathrm{SL}_2}, \mathrm{O}(2, 1))$ inside
$\widetilde{\mathrm{Sp}_4}$. The additive theta lift of
Gritsenko--Cl\'ery $2018$ (arXiv:$1802.00269$) attaches a paramodular
Siegel form of weight $k + 1/2$ on $\widetilde{\mathrm{Sp}_4}$ to a
Jacobi form of weight $k + 1/2$ and index $1$, using the Howe pair
$(\widetilde{\mathrm{Sp}_4}, \mathrm{O}(1, 1))$ inside
$\widetilde{\mathrm{Sp}_8}$.

\emph{Bezrukavnikov attack.} The first-wave spine wrote
``metaplectic cover $\Mpr_4(\mathbb{A})$'' without distinguishing
(a) the genuine twofold cover versus a splitting pullback,
(b) the local versus adelic statement of the cover, (c) the level
subgroup at each finite place. The second-wave rectification corrected
the Howe-pair attribution; the third-wave adds the level-and-cover
discipline. \emph{Heal.} Every metaplectic statement names the cover
(twofold genuine / fourfold at $N = 7$), the level subgroup, and
the local versus adelic side; half-integer weights refer to genuine
representations of the cover, not pullbacks. The ``$1/4$ at $N = 7$''
is the weight of a genuine representation of the fourfold
central extension $\widetilde{\mathrm{Sp}_4}^{(4)}$ by $\mu_4$, not of
the ordinary twofold cover (Kubota $1969$
\emph{J.~Math.~Soc.~Japan}~$21$:$263$; Matsumoto $1969$
\emph{Ann.~Sci.~\'Ec.~Norm.~Sup\'er.}~$2$:$1$ for the general
construction of higher metaplectic covers).
\end{theorem}

\subsection{Cycle 3: Motivic Galois super-Lie algebra coefficient
ring}
\label{wn:subsec:third-pass-MGSL}

\begin{theorem}[Motivic Galois super-Lie algebra: coefficient ring and
category of motives]
\label{wn:thm:third-pass-MGSL}\ClaimStatusTheorem

The motivic Galois super-Lie algebra $\mathfrak{mgsl}$ associated to
the unified BKM crown is defined in the super-Tannakian category
\[
\mathrm{MT}(\QQ, \QQ)^{\mathrm{super}}
\;=\; \text{super-mixed-Tate motives over }\QQ\text{ with }\QQ\text{
coefficients},
\]
in the sense of Deligne $1989$ / Deligne--Goncharov $2005$
(\emph{Ann.~Sci.~\'Ec.~Norm.~Sup\'er.}~$38$:$1$) on the ungraded
mixed-Tate motives, tensored with the $\ZZ/2$-supercategory of
$V^{s\natural}$-type Conway-super denominators (Duncan $2007$
\emph{Adv.~Math.}~$214$:$569$). The coefficient ring on Betti
realisation is $\QQ$; on $\ell$-adic realisation it is $\QQ_\ell$; the
comparison is via Grothendieck's motivic Betti-\'etale comparison,
with Ayoub $2014$ / Cisinski--D\'eglise $2019$
(\emph{Triangulated Categories of Mixed Motives}, Springer) supplying
the $\infty$-categorical refinement.

The super-Lie algebra $\mathfrak{mgsl}$ is
\[
\mathfrak{mgsl} \;=\; \mathrm{Lie}\bigl(\pi_1^{\mathrm{mot}}(\mathrm{MT}
(\QQ, \QQ)^{\mathrm{super}}, \omega_{\mathrm{Betti}})\bigr),
\]
the super-Lie algebra of the motivic fundamental pro-algebraic
super-group with fibre functor $\omega_{\mathrm{Betti}}$; its
generators include the odd Koszul elements $\sigma_{2n+1}$ of weight
$-2(2n+1)$ for $n \geq 1$ attached to odd zeta values, plus the
Conway-super odd generators attached to $V^{s\natural}$ moonshine
(Duncan--Mack-Ono $2015$ \emph{Res.~Math.~Sci.}~$2$:$26$).

\emph{Bezrukavnikov attack.} A statement of the form ``$\mathfrak{mgsl}$
acts on the BKM crown'' is meaningless without naming (a) the
coefficient ring ($\QQ$, $\QQ_\ell$, $\QQ_p$, or motivic), (b) the
Tannakian category of motives ($\mathrm{MM}(\QQ)$ mixed,
$\mathrm{MT}(\QQ)$ mixed-Tate, or Nori motives $\mathrm{NMM}(\QQ)$),
(c) the fibre functor selecting $\pi_1^{\mathrm{mot}}$, (d) the
super-refinement's base.

\emph{Heal.} For the programme-canonical identification of
$\mathfrak{mgsl}$ with automorphisms of the Borcherds-crown
denominator algebra, the motivic category is
$\mathrm{MT}(\QQ, \QQ)^{\mathrm{super}}$ (super-mixed-Tate motives
over $\QQ$ with $\QQ$ coefficients), the fibre functor is Betti
realisation $\omega_B\colon \mathrm{MT}(\QQ)^{\mathrm{super}}
\to \mathrm{sVect}_{\QQ}$, the pro-algebraic motivic-Galois group is
$\pi_1^{\mathrm{mot}}(\mathrm{MT}(\QQ)^{\mathrm{super}}, \omega_B)$
(Deligne--Goncharov $2005$ super-refined), and $\mathfrak{mgsl}$ is
its super-Lie algebra. Every inscription of $\mathfrak{mgsl}$ in the
monograph names this four-tuple.

\emph{$\ell$-adic realisation.} The $\ell$-adic avatar
$\mathfrak{mgsl}_\ell = \mathfrak{mgsl} \otimes_{\QQ} \QQ_\ell$ acts
on $\ell$-adic \'etale cohomology of the Borcherds-crown objects via
the $\ell$-adic realisation functor $\omega_\ell$; the comparison
isomorphism $\omega_B \otimes_\QQ \QQ_\ell \simeq \omega_\ell$ is
part of the motivic Betti-\'etale comparison (Grothendieck's standard
conjecture; unconditional on $\mathrm{MT}(\QQ)$ by Deligne).
\end{theorem}

\begin{theorem}[$\mathrm{GRT}_1$ action, compatible with which
formulation]
\label{wn:thm:third-pass-GRT1}\ClaimStatusConjectured

The Grothendieck--Teichm\"uller Lie algebra $\mathfrak{grt}_1$ acts on
$\mathfrak{mgsl}$ via the Drinfeld associator torsor
(Drinfeld $1990$ \emph{Leningrad Math.~J.}~$2$:$829$, Thm.~A$'$);
compatibly in the sense of Ihara $1991$ (\emph{Proc.~ICM Kyoto}
p.~$99$) with the Ihara action on the motivic pro-unipotent
fundamental group of $\bP^1_{\QQ} \setminus \{0, 1, \infty\}$.
Under Furusho $2010$ (\emph{Ann.~Math.}~$171$:$545$) the
Drinfeld--Kohno and Ihara formulations coincide; the monograph uses
this identified action.

\emph{$\mathrm{GRT}_1$-transitivity on the full BKM crown is
conditional and open.} The claim ``$\mathrm{GRT}_1$ acts transitively
on the full BKM crown'' requires (a) a precise target on which
transitivity is asserted (the set of primary lifts, the deformation
space of the motivic Galois super-Lie algebra, or the moduli of
Borcherds products); (b) the generic choice of associator in
$\mathrm{GRT}_1(\QQ)$ that enacts the transitivity; (c) the
compatibility condition between the Drinfeld--Kohno and Ihara actions
on the motivic-Galois target. Under the strongest interpretation --
$\mathrm{GRT}_1(\QQ)$ acts transitively on the set of $\QQ$-rational
points of the deformation space of $\mathfrak{mgsl}$-module structures
on the BKM crown -- the statement is open. Primary evidence comes
from the Alekseev--Enriquez--Torossian $2010$ framework and the
Willwacher $2015$ graph-complex realisation of $\mathfrak{grt}_1$
(\emph{Invent.~Math.}~$200$:$671$).

\emph{Bezrukavnikov attack.} A $\mathrm{GRT}_1$-transitivity statement
that does not name the formulation (Drinfeld--Kohno associator versus
Ihara motivic versus Willwacher graph-complex) and the transitivity
target is not a theorem. \emph{Heal.} The monograph-canonical
$\mathrm{GRT}_1$ refers to the pro-algebraic $\QQ$-group of
Drinfeld--Kohno associators (Drinfeld $1990$ Thm.~A$'$), identified
with Ihara's motivic $\mathrm{GRT}_1$ via Furusho $2010$, with
transitivity claim on the $\QQ$-rational deformation-space of
$\mathfrak{mgsl}$-module structures on the BKM crown left open at
this level of primary-source support. Willwacher's graph-complex
isomorphism $\mathfrak{grt}_1 \simeq H^0(\mathrm{GC}_2)$ allows the
computation to be performed in the graph-complex presentation, with
agreement deduced from Furusho.
\end{theorem}

\subsection{Cycle 4: Chenevier determinant versus pseudo-character}
\label{wn:subsec:third-pass-Chenevier}

\begin{theorem}[Chenevier determinants on the paramodular Galois
representations]
\label{wn:thm:third-pass-Chenevier}\ClaimStatusTheorem

The Galois-representation invariants attached to paramodular Siegel
cusp forms $F \in S_k(K(t))$ (via the Weissauer $2005$
\emph{Compos.~Math.}~$141$:$127$ / Laumon $2005$
\emph{Ast\'erisque}~$302$:$1$ correspondence for $F$ of degree~$2$
spinor $L$-function) are encoded in the monograph by the Chenevier
determinant, not by the Taylor--Rouquier pseudo-character. Following
Chenevier $2014$ (\emph{Algebra Number Theory}~$8$:$1527$,
Thm.~$2.22$), a $d$-dimensional determinant on a profinite group
$G$ valued in a commutative ring $A$ is an $A$-valued polynomial
$d$-law $D\colon A[G] \to A$ of degree $d$, satisfying $D(1) = 1$
and multiplicativity $D(xy) = D(x)D(y)$ for $x, y \in A[G]$. For
$A = A_{\mathfrak{m}}$ a local Noetherian complete $\ZZ_p$-algebra
with residue field $k$ and $G = \mathrm{Gal}(\overline{\QQ}/\QQ)$,
a $d$-dimensional Chenevier determinant $D$ with residual
representation
$\overline{\rho}\colon G \to \mathrm{GL}_d(k)$ factors through a
unique $A$-algebra morphism
$R_{\overline{\rho}}^{\mathrm{det}} \to A$ from the universal
deformation ring of determinants lifting $\overline{\rho}$ (Chenevier
$2014$ Thm.~$3.1$).

For the paramodular $\Delta_5$ Galois representation
$\rho_{\Delta_5, \ell}\colon G_\QQ \to \mathrm{GSp}_4(\QQ_\ell)$
(conjecturally; Calegari--Geraghty $2018$
\emph{Invent.~Math.}~$211$:$297$ for automorphy of certain
paramodular forms), the associated Chenevier determinant
$D_{\Delta_5} = \det \circ \rho_{\Delta_5, \ell}$ carries the
spin-$L$-function information in the coefficient sequence
$\{\lambda_p(\Delta_5)\}$.

\emph{Why Chenevier determinants, not pseudo-characters.} The
Taylor--Rouquier pseudo-character
$T\colon G \to A$ (Taylor $1991$ \emph{Invent.~Math.}~$103$:$289$;
Rouquier $1996$ \emph{J.~Algebra}~$180$:$571$) is defined only in
characteristic $p > d$; in small residue characteristic ($p \leq d$),
the pseudo-character loses the determinant data and fails to recover
the representation. Chenevier determinants remedy this: they work in
arbitrary residue characteristic, including $p \leq d$, and recover
the semisimple representation uniquely on the algebraic closure
(Chenevier $2014$ Thm.~$2.12$). For paramodular forms of weight
$k \geq 3$ and degree $2$ spinor representations, $d = 4$; residue
characteristic $p = 2, 3$ falls into the pseudo-character-defective
range, forcing the monograph to use Chenevier determinants.

\emph{Bezrukavnikov attack.} The first-wave spine invoked ``the
pseudo-character of $\rho_{\Delta_5}$'' without naming (a) the
residue characteristic, (b) whether the pseudo-character is defined
in that characteristic, (c) which residual representation
$\overline{\rho}$ is selected, (d) whether the pseudo-character
versus the Chenevier determinant is used.
\emph{Heal.} Every Galois-representation invariant in the monograph
is a Chenevier determinant; the universal deformation ring is
$R_{\overline{\rho}}^{\mathrm{det}}$; the residual representation
$\overline{\rho}\colon G_\QQ \to \mathrm{GSp}_4(\mathbb{F}_\ell)$ is
named at every use with its Artin conductor and inertia type. The
pseudo-character terminology is reserved for the special case
$p > 4$ where Taylor--Rouquier applies and the two notions agree.
\end{theorem}

\subsection{Cycle 5: Humbert divisor discriminant conventions}
\label{wn:subsec:third-pass-Humbert}

\begin{theorem}[Humbert divisor $H_n$ in $\mathcal{A}_2$: discriminant
convention fixed]
\label{wn:thm:third-pass-Humbert}\ClaimStatusTheorem

The Humbert surfaces $H_n \subset \mathcal{A}_2(\CC)$ of discriminant
$n \in \ZZ_{\geq 1}$ are the irreducible codimension-one subvarieties
of the Siegel modular threefold $\mathcal{A}_2 =
\mathrm{Sp}_4(\ZZ)\backslash\mathbb{H}_2$ parametrising abelian
surfaces $A$ with real multiplication by the order of discriminant
$n$ in a real quadratic field $F = \QQ(\sqrt{n})$ (or by
$\ZZ \times \ZZ$ at $n = 1$ split case), equivalently the loci of
Jacobians of curves of genus $2$ admitting an isogeny of that
discriminant. Primary: Humbert $1899$ (\emph{J.~Math.~Pures
Appl.}~$5$:$129$); van~der~Geer $1988$ \emph{Hilbert Modular
Surfaces} Ch.~IX; Gritsenko--Hulek $1998$ \emph{Abelian Varieties}
\S$1.5$.

\emph{Normalisation convention.} The monograph adopts the
van~der~Geer convention: $H_n$ is irreducible for $n$
squarefree, and
\[
H_n = \bigl\{\,[A] \in \mathcal{A}_2 \;\bigm|\;
\mathrm{End}(A)\otimes\QQ \supseteq \QQ(\sqrt{n})\,\bigr\},
\qquad
\mathrm{codim}_{\mathcal{A}_2}(H_n) = 1.
\]
At $n = 1$: $H_1$ parametrises products of elliptic curves
$A = E_1 \times E_2$ with the split real-multiplication locus; the
diagonal locus $\{E \times E : E \text{ elliptic}\}$ is a
subvariety of $H_1$ of codimension $1$ in $H_1$ (hence codim $2$ in
$\mathcal{A}_2$), not all of $H_1$. The Humbert discriminant
convention here is $n = $ discriminant of the order of real
multiplication.

\emph{Bezrukavnikov attack.} The first-wave spine wrote
``Humbert divisor $H_1 \subset \mathcal{A}_2$'' without disambiguating
(a) Humbert convention used (van~der~Geer discriminant, Gritsenko--
Hulek discriminant, or Shimura discriminant -- these differ by
factors of $4$ at non-fundamental discriminants), (b) whether $H_1$
means the full locus of products-of-elliptic-curves or the diagonal
elliptic-product locus (codim $1$ versus codim $2$), (c) whether $H_4$
(``fourfold-discriminant'') invoked in the Vol~III Humbert-tower
sibling discussion uses $n = 4$ (non-fundamental, equal to Humbert
$H_{\mathrm{disc = 4}}$) or $d = 2$ in the Gritsenko fundamental-
discriminant convention.

\emph{Heal.} The monograph uses van~der~Geer $1988$ Ch.~IX
discriminant convention throughout: $H_n$ is the Humbert surface of
discriminant $n \geq 1$ parametrising abelian surfaces with real
multiplication by the order of discriminant $n$. The residue
formula
\[
\mathrm{Res}_{H_1}\bigl[(2\pi i z)^2 \Delta_5^{-2}\bigr]\Big|_{z = 0}
\;=\; \frac{1}{\eta(\tau)^{24}\,\eta(\sigma)^{24}},
\]
cited via Gritsenko--Nikulin $1996$ Cor.~$5.2$ and
Gritsenko--Hulek $1997$, refers to the Humbert divisor $H_1$ in the
van~der~Geer convention, not to the codim-$2$ diagonal. The
diagonal-specialisation $\tau = \sigma$ performs a second restriction
from $H_1$ to the codim-$2$ diagonal, producing $\eta(q)^{-48}$. Every
Humbert-divisor invocation in the monograph specifies the convention
(van~der~Geer, default) and the residue specialisation (full-$H_n$
versus diagonal-$\tau = \sigma$).

\emph{Modular-curve companion.} The $N$-CHL twined family
$\{H_N\}_{N \in \{1, 2, 3, 4, 6\}}$ of
Theorem~\ref{wn:thm:second-pass-single-ladder} is compatible with the
Humbert-discriminant convention via the identity
$\mathrm{disc}(\ZZ[\sqrt{N^2}]) = 4N^2$ for $N$ squarefree, mapping
$H_N$-twined paramodular data to $H_{N^2}$-discriminant Humbert data
up to the non-fundamental-discriminant correction. The two families
are conceptually distinct even when their labels agree numerically.
\end{theorem}

\subsection{Cycle 6: Class-$\mathcal{S}$ gauge and flavour content,
fully named}
\label{wn:subsec:third-pass-class-S}

\begin{theorem}[Class-$\mathcal{S}$ $(A_1, \Sigma_{0, 24})$: gauge
group and flavour structure inscribed]
\label{wn:thm:third-pass-class-S-gauge-flavour}\ClaimStatusTheorem

The four-dimensional superconformal field theory
$\mathcal{T}[A_1, \Sigma_{0, 24}]$ is the class-$\mathcal{S}$ theory
of Gaiotto $2012$ (\emph{J.~High Energy Phys.}~$2012$:$34$)
associated to the pair:
\begin{itemize}
\item \emph{ADE label} $A_1$, equivalently the
$6$d $\mathcal{N} = (2, 0)$ parent theory of type $\mathfrak{su}(2)$,
whose tensor-branch coefficients are the simple roots of
$\mathfrak{su}(2) = A_1$;
\item \emph{Riemann surface data} $\Sigma_{0, 24}$: a
$24$-punctured sphere with full maximal regular punctures (the maximal
nilpotent orbit in the $A_1$ Lie algebra) at each puncture, equivalently
every puncture carries $\mathrm{SU}(2)$-flavour symmetry;
\item \emph{Gauge group after a pants decomposition}:
$\mathrm{SU}(2)^{n - 3} = \mathrm{SU}(2)^{21}$ from internal tubes;
the $22$ trinions carry half-hypermultiplets in the pseudoreal
$(\mathbf{2}, \mathbf{2}, \mathbf{2})$ representation of
$\mathrm{SU}(2)^3$;
\item \emph{Flavour symmetry}: $\mathrm{SU}(2)^{24}$ from the external
punctures, producing $(n_v, n_h) = (63, 88)$;
\item \emph{Coulomb-branch dimension}: $n - 3 = 21$ from the
$\mathrm{SU}(2)^{21}$ gauge sector.
\end{itemize}
The central charges follow from Shapere--Tachikawa $2008$
(\emph{J.~High Energy Phys.}~$2008$:$109$) trace-anomaly formula
$c_{4d} = (2n_v + n_h)/12 = (2 \cdot 63 + 88)/12 = 214/12 = 107/6$.

\emph{Bezrukavnikov attack.} The first-wave spine wrote
``$\mathcal{T}[A_1, \Sigma_{0, 24}]$ with $c_{4d} = 107/6$'' without
naming (a) the $6$d parent ADE type ($A_1$ versus $A_{N - 1}$ for
$N \geq 3$, which produce different theories with different central
charges), (b) the puncture type at each external puncture (maximal
regular versus minimal nilpotent versus irregular), (c) the gauge
group resulting from a specific pants decomposition, (d) the flavour
symmetry associated to the puncture choice. Without these, the
class-$\mathcal{S}$ theory is not specified; different puncture
choices produce different SCFTs with the same $6$d parent.

\emph{Heal.} The monograph-canonical class-$\mathcal{S}$ theory at
$\Sigma_{0, 24}$ is $\mathcal{T}[A_1, \Sigma_{0, 24}, \{\text{maximal
regular punctures}\}]$, with gauge group $\mathrm{SU}(2)^{21}$ after
a generic pants decomposition, matter half-hypermultiplets in the
trinion pseudoreal representation, and flavour symmetry
$\mathrm{SU}(2)^{24}$. The $2$d chiral-algebra map of Beem--Rastelli
$2018$ (\emph{J.~High Energy Phys.}~$2018$:$114$) sends this SCFT to
the vertex algebra $\mathcal{V}_{4d}[\mathcal{T}[A_1, \Sigma_{0, 24}]]$
with $c_{2d} = -12 c_{4d} = -214$. Every class-$\mathcal{S}$
invocation names the ADE type, the punctured Riemann surface, the
puncture profile, the gauge group after pants, and the matter / flavour
content.
\end{theorem}

\subsection{Cycle 7: $8\times 5$ bistrata categorical framework}
\label{wn:subsec:third-pass-bistrata}

\begin{theorem}[The $8\times 5$ bistrata matrix: $(\infty, 1)$-
categorical framework]
\label{wn:thm:third-pass-bistrata}\ClaimStatusTheorem

The ``$8 \times 5$ bistrata matrix'' organises
$8$ ambient-geometric strata against $5$ archetype-algebraic columns
(Heisenberg / affine Kac--Moody / $\beta\gamma$ / Virasoro / Mukai-
enhanced Heisenberg), indexing $40$ inscription cells in the product
$(\text{ambient}) \times (\text{archetype})$. The precise
$(\infty, 1)$-categorical framework is:
\[
\mathcal{B} \;=\;
\mathrm{Fun}\bigl(\mathcal{A}^{\mathrm{op}} \times \mathcal{C},\;
\mathrm{ChirAlg}^{E_1}_{\mathrm{ch}}\bigr)
\]
where:
\begin{itemize}
\item $\mathcal{A}$ is the $(\infty, 1)$-category of ambient
configurations with objects
$\{\mathrm{pt},\RR, \CC, \HH, \mathbb{D}, A^2\setminus\mathrm{pt},
\Sigma_g, \Sigma_g \times_{\mathrm{top}} E\}$ ($8$ objects) and
morphisms given by Whitehead-tower compatibility maps of framed
manifolds with tubular neighbourhoods (Ayala--Francis $2015$
\emph{Geom.~Topol.}~$19$:$2463$; Lurie \emph{Higher Algebra}~$5.5$);
\item $\mathcal{C}$ is the set $\{G, L, C, M, B\}$ of five archetypes,
regarded as discrete $(\infty, 1)$-category with identity morphisms
only;
\item $\mathrm{ChirAlg}^{E_1}_{\mathrm{ch}}$ is the presentable
stable $(\infty, 1)$-category of $E_1$-chiral algebras on a smooth
curve $C$ over $\CC$ in the sense of Beilinson--Drinfeld $2004$
Ch.~$3.4$ / Francis--Gaitsgory $2012$.
\end{itemize}

The monograph's bistrata cells are the evaluation values of objects
$F \in \mathcal{B}$ on pairs $(a, c) \in \mathcal{A}^{\mathrm{op}}
\times \mathcal{C}$. Because $\mathcal{A}$ is not discrete, the
bistrata matrix is the underlying object set of the functor
category; morphisms between cells in the same column correspond to
ambient-topology changes (e.g., $\mathrm{pt} \to \RR \to \CC \to \HH$
corresponds to the dimensional-extension tower).

\emph{Bezrukavnikov attack.} The first-wave spine presented the
$8 \times 5$ matrix as a raw $40$-cell table without (a) specifying
the $(\infty, 1)$-category of ambient configurations, (b) indicating
that columns are discrete and rows have non-trivial morphisms, (c)
specifying whether the cells are $E_1$-chiral algebras or factorisation
algebras, (d) naming the base curve over which the chiral algebras
live. \emph{Heal.} The bistrata is the $\mathrm{Ob}$-truncation of the
functor $(\infty, 1)$-category $\mathcal{B}$ described above; columns
are indexed by the discrete archetype set $\{G, L, C, M, B\}$;
rows are indexed by objects of the non-discrete Ayala--Francis
ambient-geometry $(\infty, 1)$-category; values are $E_1$-chiral
algebras on a fixed reference curve $C$ over $\CC$. Stage-$2$
specialisation $\mathrm{Sp}^{\mathrm{ch}}_{\Sigma_{d-1}, C}$ realises
the row-indexed ambient-change morphisms as functors between chiral
algebras; Theorem~A (bar--cobar adjunction, Vol~I) constrains the
column structure via the four-archetype dichotomy plus the
Mukai-enhanced $\mathsf{B}$-row five-archetype extension.
\end{theorem}

\subsection{Third-pass residual frontier}
\label{wn:subsec:third-pass-frontier}

\begin{theorem}[Third-pass frontier items]
\label{wn:thm:third-pass-frontier}

The Bezrukavnikov-precision pass closes several third-pass gaps and
opens others. Closed: (i) D-module categories named with
regular-holonomicity and left-$\mathcal{D}$-module choice;
(ii) Shimura variety of Hodge type for Kuga--Satake, with reductive
group $\mathrm{GSpin}(P^2)$ and level $K$ named; (iii) Chenevier
determinants with residual-representation data replace
pseudo-characters throughout; (iv) Humbert-discriminant convention
fixed to van~der~Geer; (v) class-$\mathcal{S}$ gauge and flavour
content inscribed. Opened:
\begin{itemize}
\item The $\mathrm{GRT}_1$-transitivity of the motivic-Galois action
on the BKM crown: open as a transitivity statement on the
deformation space of $\mathfrak{mgsl}$-module structures;
Willwacher $2015$ graph-complex presentation provides a concrete
target for verification.
\item Integrality of Chenevier determinant lifts of paramodular
Galois representations at residue characteristic $p = 2, 3$: open
for $\Delta_5$ in the absence of a Calegari--Geraghty-type
automorphy theorem at these primes.
\item $(\infty, 1)$-functoriality of the Ayala--Francis ambient
$(\infty, 1)$-category morphisms with the Stage-$2$ specialisation:
the row-morphism / Stage-$2$-functor compatibility requires a
strictification of the Ayala--Francis presentation compatible with
Drinfeld--Gaitsgory $2015$'s factorisation-sheaf framework.
\item Compatibility of the Saito $\mathrm{MHM}$ six-functor formalism
with the factorisation-algebra $E_3^{\mathrm{hol}}$-structure on
compact CY$_3$: the relative decomposition for the
product-$K3\times E$ proper map is Saito $1990$ Thm.~$0.1$, but its
integration with the Costello--Gwilliam $\cF_X$ factorisation
structure requires a translation between mixed Hodge modules and
the $\mathcal{D}$-module lane of Beilinson--Drinfeld.
\end{itemize}
\end{theorem}

\subsection{Third-pass synthesis, in one sentence}
\label{wn:subsec:third-pass-one-line}

\emph{The geometric-representation-theoretic substrate of the K3
$\times$ E crystal is organised by four data -- D-module category
(regular-holonomic left $\mathcal{D}$-modules on the Ran space over
$\CC$), Shimura variety type (Hodge-type $\mathrm{Sh}_K(\mathrm{GSpin}
(P^2))$ with Clifford embedding into Siegel $\mathcal{A}_g$),
motivic-Galois coefficient ring (super-mixed-Tate motives over $\QQ$
with $\QQ$ coefficients), and Chenevier determinant invariants
(with named residual representation $\overline{\rho}_{\Delta_5}$ over
$\mathbb{F}_\ell$) -- and is compatible with the second-pass
four-functor square via the six-functor formalism on $\mathrm{MHM}$
(Saito), with Humbert-divisor labels in the van~der~Geer
discriminant convention, class-$\mathcal{S}$ gauge and flavour
content inscribed ($A_1$ parent, maximal regular punctures,
$\mathrm{SU}(2)^{21}$ gauge, $\mathrm{SU}(2)^{24}$ flavour,
$(n_v, n_h) = (63, 88)$), and the $8\times 5$ bistrata realised as
the object-truncation of the Ayala--Francis ambient-geometry $(\infty,
1)$-functor category $\mathrm{Fun}(\mathcal{A}^{\mathrm{op}} \times
\mathcal{C}, \mathrm{ChirAlg}^{E_1}_{\mathrm{ch}})$; every statement
in the monograph now carries its category, its hypotheses, and its
primary-source anchor at the precision where Beilinson--Bernstein,
Saito, Kashiwara, Chenevier, van~der~Geer, and Drinfeld--Gaitsgory
can each be cited for exactly the use made of them.}

% End of third-pass Bezrukavnikov refinement.
